% =============================================================================
% SRS v1.0.0 -- Sistema de Gestion de Pre-Sustentaciones UTEQ
% ISO/IEC/IEEE 29148:2018
% Compilar con: latexmk -pdf -interaction=nonstopmode -halt-on-error SRS-v1.0.0.tex
% =============================================================================
\documentclass[11pt,a4paper]{article}
\usepackage[utf8]{inputenc}
\usepackage[T1]{fontenc}
\usepackage[spanish]{babel}
\usepackage{geometry}
\geometry{letterpaper, margin=1in}
\usepackage{setspace}
\setstretch{1.15}
\usepackage{graphicx}
\usepackage{booktabs}
\usepackage{float}
\usepackage{longtable}
\usepackage{tabularx}
\usepackage{array}
\usepackage[table]{xcolor}
\usepackage{enumitem}
\usepackage{seqsplit}
\usepackage{fancyhdr}
\usepackage{hyperref}
\hypersetup{colorlinks=true, linkcolor=blue, urlcolor=blue, citecolor=blue,
            pdftitle={SRS v1.0.0 - Sistema de Gestion de Pre-Sustentaciones UTEQ}}

\pagestyle{fancy}
\fancyhf{}
\fancyhead[L]{\small SRS v1.0.0 --- Sistema de Gestión de Pre-Sustentaciones UTEQ}
\fancyhead[R]{\small\thepage}
\renewcommand{\headrulewidth}{0.4pt}

% ---- Colores de estado -------------------------------------------------------
\definecolor{verdeok}{RGB}{0,120,60}
\definecolor{ambar}{RGB}{176,110,0}
\definecolor{rojoko}{RGB}{178,34,34}
\definecolor{grisln}{RGB}{120,120,120}
\definecolor{filagris}{RGB}{242,242,242}

% ---- Macros de requisito -----------------------------------------------------
% Rutas y nombres de clase: \cd{} permite el quiebre de línea dentro del
% fragmento monoespaciado. Sin esto, un endpoint largo produce un overfull hbox
% que saca texto normativo fuera del margen derecho en el PDF.
\DeclareRobustCommand{\cd}[1]{\texttt{\protect\seqsplit{#1}}}

\newcommand{\estadoV}{\textcolor{verdeok}{\textbf{Verificado}}}
\newcommand{\estadoI}{\textcolor{ambar}{\textbf{Implementado}}}
\newcommand{\estadoP}{\textcolor{grisln}{\textbf{Planificado}}}
\newcommand{\estadoN}{\textcolor{rojoko}{\textbf{No cumplido}}}

% Encabezado de requisito. No se añade al índice: con 87 requisitos el índice
% dejaría de ser navegable. El índice llega al nivel de subsección temática.
\newcommand{\REQ}[2]{%
  \par\vspace{3mm}%
  \noindent\rule{\linewidth}{0.7pt}\par\nobreak\vspace{0.6mm}%
  \noindent{\normalsize\textbf{#1\quad #2}}\par\nobreak\vspace{0.6mm}%
  \noindent\rule{\linewidth}{0.3pt}\par\nobreak\vspace{1mm}}

\newcommand{\campo}[2]{\noindent\textbf{#1:} #2\par\vspace{0.4mm}}

\setlist[itemize]{nosep,leftmargin=1.4em,topsep=1mm}
\setlist[enumerate]{nosep,leftmargin=1.8em,topsep=1mm}

\title{}\date{}
\begin{document}

% ============================== CARÁTULA =====================================
\begin{titlepage}
\centering
\vspace*{0.8cm}
{\LARGE \textbf{UNIVERSIDAD TÉCNICA ESTATAL DE QUEVEDO}}\\[0.3cm]
{\large Facultad de Ciencias de la Ingeniería}\\[0.2cm]
{\large Carrera de Ingeniería de Software (Rediseño)}\\[0.8cm]

{\large \textbf{ASIGNATURA: Aplicaciones Web}}\\[0.3cm]
{\large Quinto Nivel --- Período Académico 2026-2027 PPA}\\[1.0cm]

{\Large \textbf{ESPECIFICACIÓN DE REQUISITOS DE SOFTWARE (SRS)}}\\[0.3cm]
{\large Estándar ISO/IEC/IEEE 29148:2018}\\[0.6cm]
{\LARGE \textbf{Sistema de Gestión de Pre-Sustentaciones UTEQ}}\\[0.6cm]
{\Large \textbf{Versión 1.0.0}}\\[0.9cm]

\small\begin{tabular}{ll}
\textbf{Repositorio:} & \url{https://github.com/carla22072004/PFC-Presustentaciones-2026} \\
\textbf{Punto del historial:} & etiqueta \texttt{v1.0.1}, \emph{commit} \texttt{9d7abbd} \\
\textbf{Rama principal al cierre:} & \emph{commit} \texttt{73c8902} (2026-09-07) \\
\textbf{DOI del software:} & \href{https://doi.org/10.5281/zenodo.22445216}{10.5281/zenodo.22445216} \\
\end{tabular}\normalsize\\[1.0cm]

{\large \textbf{EQUIPO DE DESARROLLO:}}\\[0.3cm]
{\large Alava Alvarado, Jean Pierre \quad (\href{https://orcid.org/0009-0001-2878-2919}{ORCID: 0009-0001-2878-2919})}\\[0.1cm]
{\large Moncayo Loor, Xavier Alejandro}\\[0.1cm]
{\large Zamora Arias, Carla Esthefania}\\[0.1cm]
{\large Barreto Rosado, Heider Dominick}\\[0.8cm]

{\large \textbf{DOCENTE-DIRECTOR:}}\\[0.3cm]
{\large Ing. Guerrero Ulloa Gleiston Cíceron, Mg.}\\[0.8cm]

{\large \textbf{FECHA DE EMISIÓN:} 9 de septiembre de 2026}\\[0.3cm]
\end{titlepage}

\tableofcontents
\newpage

% ====================== 1. IDENTIFICACIÓN DEL DOCUMENTO ======================
\section{Identificación y control del documento}

\subsection{Historial de versiones}

\begin{longtable}{p{1.6cm} p{1.8cm} p{2.6cm} p{8.2cm}}
\toprule
\textbf{Versión} & \textbf{Fecha} & \textbf{Autor} & \textbf{Cambios} \\
\midrule
\endfirsthead
\toprule
\textbf{Versión} & \textbf{Fecha} & \textbf{Autor} & \textbf{Cambios} \\
\midrule
\endhead
1.0.0 & 2026-09-09 & Equipo de desarrollo &
\textbf{Primera versión de esta especificación.} Corpus de \textbf{64 requisitos funcionales} y
\textbf{26 no funcionales}, redactados y verificados contra el código real de la etiqueta
\texttt{v1.0.1}. Incorpora la plantilla completa por requisito (enunciado normativo,
\emph{rationale}, criterio de aceptación, método de verificación y estado) y las cinco secciones
estructurales que la norma exige: interfaces externas, máquinas de estado del dominio, matriz de
permisos rol\,$\times$\,operación, correspondencia con el Anexo C y cláusula de verificación. \\
\bottomrule
\caption{Historial de versiones del SRS. La bitácora completa de cambios de requisitos
vive en \cd{docs/requisitos/CHANGELOG-REQ.md}.}
\end{longtable}

\noindent\textbf{Documentos a los que reemplaza.} Esta especificación sustituye a los dos
borradores de requisitos anteriores del proyecto, que quedan archivados en
\cd{docs/requisitos/historico/} y no deben citarse como especificación vigente: el borrador del
2026-07-30 (5 historias de usuario, sin casos de uso) y el del 2026-08-17 (12 historias y 12 casos
de uso, con un resumen de 12 requisitos funcionales y 4 no funcionales). Ninguno de los dos
describía más de una novena parte de las rutas que el sistema expone hoy; las correcciones
concretas que se les hicieron están en la sección~\ref{sec:cambios}.

\subsection{Alcance de esta versión y punto de corte}

Todo enunciado de estado, cifra y referencia a código de este documento se verificó
\textbf{contra el árbol de archivos de la etiqueta \texttt{v1.0.1}} (\emph{commit}
\texttt{9d7abbd}, 2026-09-05) y se contrastó con la rama principal al \emph{commit}
\texttt{73c8902} (2026-09-07). Las cifras que dependen de una medición traen la fecha de esa
medición dentro del propio requisito: una cifra sin fecha envejece sola y deja de ser
verificable.

\subsection{Autoridad de cada artefacto}
\label{sec:autoridad}

Este documento y la matriz de trazabilidad no compiten: cada uno es fuente de verdad de algo
distinto, y esa separación es lo que evita que se desincronicen sin que nadie lo note.

\begin{table}[H]\centering\small
\begin{tabular}{p{4.2cm} p{10.5cm}}
\toprule
\textbf{Artefacto} & \textbf{Es fuente de verdad de} \\
\midrule
Este SRS & El enunciado normativo, el \emph{rationale}, la prioridad MoSCoW, el criterio de
aceptación y el método de verificación de cada requisito. \\
\cd{docs/trazabilidad/matriz.csv} & El estado de verificación, el endpoint, el módulo, la
clase de prueba y la evidencia empírica de cada requisito. \\
\cd{scripts/validate-traceability.sh} & La coherencia entre ambos y con el disco. Ninguna
versión se publica si el validador no termina con código 0. \\
\bottomrule
\end{tabular}
\end{table}

% ========================== 2. INTRODUCCIÓN ==================================
\section{Introducción}

\subsection{Propósito}
Este documento define los requisitos funcionales y no funcionales del Sistema de Gestión de
Pre-Sustentaciones UTEQ conforme al estándar ISO/IEC/IEEE 29148:2018. Sirve como contrato
técnico entre la coordinación académica, los docentes, los estudiantes y el equipo
desarrollador, y como base de verificación: cada requisito declara cómo se comprueba y en qué
estado está.

\subsection{Alcance del sistema}
El sistema automatiza el ciclo completo de pre-sustentación de trabajos de titulación:
provisión de cuentas y control de acceso por permisos; registro y revisión de la solicitud;
asignación de tutor y acompañamiento en tres fases de tutoría con entrega de correcciones;
carga del anteproyecto en PDF con verificación de integridad SHA-256; conformación del
tribunal; programación del cronograma de defensa con detección de conflictos de sala y de
jurado; evaluación por rúbrica y cálculo de la nota final ponderada; emisión del acta en PDF
y su firma digital multi-actor; notificaciones; reportes de gestión; un centro de orientación
de titulación; una bitácora de auditoría; y la gestión de respaldos y recuperación de la base
de datos.

\textbf{Fuera de alcance declarado:} matrícula y expediente académico institucional, pago de
aranceles, firma electrónica con validez jurídica ante entidad certificadora, y la sustentación
final de grado (este sistema cubre la \emph{pre}-sustentación).

\subsection{Definiciones, acrónimos y abreviaturas}
\begin{itemize}
    \item \textbf{SRS:} \emph{Software Requirements Specification} (ISO/IEC/IEEE 29148:2018).
    \item \textbf{UTEQ:} Universidad Técnica Estatal de Quevedo.
    \item \textbf{RF / RNF:} Requisito Funcional / Requisito No Funcional.
    \item \textbf{HU / CU:} Historia de Usuario (Connextra) / Caso de Uso (plantilla Cockburn).
    \item \textbf{MoSCoW:} priorización \emph{Must} / \emph{Should} / \emph{Could} / \emph{Won't}.
    \item \textbf{JWT:} \emph{JSON Web Token}. \textbf{jti}: identificador único del token.
    \item \textbf{SP:} procedimiento almacenado PL/pgSQL.
    \item \textbf{IDOR:} \emph{Insecure Direct Object Reference}, acceso a un recurso ajeno por
          manipulación de su identificador.
    \item \textbf{RPO / RTO:} objetivo de punto de recuperación / objetivo de tiempo de
          recuperación.
    \item \textbf{Titular:} persona a la que pertenecen los datos personales tratados.
    \item \textbf{Tribunal:} conjunto de tres docentes con roles \texttt{PRESIDENTE},
          \cd{VOCAL\_1} y \cd{VOCAL\_2} asignados a una solicitud.
\end{itemize}

\subsection{Convención de redacción normativa}
\label{sec:convencion}

Esta convención es de aplicación obligatoria a todo requisito de este documento y a todo
requisito que se agregue en el futuro.

\begin{enumerate}
    \item El verbo vinculante es \textbf{«debe»}. Todo enunciado que use «debe» expresa una
          obligación exigible y verificable.
    \item \textbf{No se usa «debería»} en el enunciado de ningún requisito. En la convención de
          la norma «debería» expresa una recomendación no vinculante y, por lo tanto, no enuncia
          un requisito. Una recomendación va en el \emph{rationale} o en una nota, nunca en el
          enunciado.
    \item \textbf{No se usa el presente indicativo descriptivo} («el sistema usa PostgreSQL»)
          para enunciar un requisito: eso describe lo construido en vez de especificar lo
          exigido. Las tecnologías impuestas se declaran como \textbf{restricciones de diseño}
          (sección~\ref{sec:restricciones}), no como requisitos.
    \item Todo requisito es \textbf{singular}: enuncia una sola capacidad. Cuando por cohesión
          se agrupan varias capacidades bajo un mismo identificador, el requisito trae
          \textbf{un criterio de aceptación por cada capacidad agrupada}; si sus estados pueden
          diferir, el requisito se divide en sufijos (\texttt{a}, \texttt{b}, \texttt{c}).
    \item El criterio de aceptación se expresa como \textbf{comportamiento observable por el
          cliente} (código HTTP y contenido de la respuesta), nunca como detalle de
          implementación interna. El nombre de la excepción o de la clase Java pertenece a la
          matriz de trazabilidad, no al criterio.
    \item El requisito \textbf{no cita rutas de archivo ni nombres de clase}: son lo primero que
          envejece con un renombrado de paquetes. El requisito cita el \textbf{endpoint}, que es
          estable; la ruta y la clase viven en la matriz, donde el validador las comprueba
          contra el disco.
\end{enumerate}

\subsection{Vocabulario cerrado de estado del requisito}
\label{sec:estados-req}

El campo \textbf{Estado} de cada requisito admite exactamente estos cuatro valores y ningún
otro. Cualquier matiz («implementado pero sin resolución ética», «medido el 18-08») va en el
campo \textbf{Nota}, nunca en el campo Estado; de lo contrario la matriz deja de poder
filtrarse y contarse, y el vocabulario se degrada en cada ampliación.

\begin{table}[H]\centering\small
\begin{tabular}{p{3.2cm} p{11.5cm}}
\toprule
\textbf{Valor} & \textbf{Significado técnico} \\
\midrule
\estadoV & La capacidad está construida \textbf{y} existe una prueba automatizada o una
evidencia empírica archivada y fechada que comprueba su criterio de aceptación. \\
\estadoI & La capacidad está construida y responde en el endpoint declarado, pero su criterio
de aceptación aún no tiene prueba ni evidencia archivada. \\
\estadoP & La capacidad \textbf{no está construida}. Existe el requisito, no la funcionalidad. \\
\estadoN & La capacidad está construida pero \textbf{no alcanza el umbral} que su propio
criterio de aceptación exige. Es distinto de «pendiente de evidencia»: aquí sí se midió, y el
resultado no cumple. \\
\bottomrule
\end{tabular}
\end{table}

\noindent\fbox{\parbox{0.98\textwidth}{\textbf{Regla de lectura:} ningún requisito marcado
\estadoP{} o \estadoN{} debe interpretarse como funcionalidad entregada y conforme. Un
requisito \emph{Must} en estado \estadoP{} o \estadoN{} es deuda declarada, y aparece como tal
en la sección~\ref{sec:cumplimiento}.}}

\subsection{Métodos de verificación}
\label{sec:metodos}

Todo requisito declara su método de verificación con este vocabulario cerrado de cuatro
valores. Declarar el método, y no solo decir «se probó», es lo que permite saber si un
requisito que no se puede cubrir con una prueba unitaria está o no verificado.

\begin{table}[H]\centering\small
\begin{tabular}{p{2.9cm} p{11.8cm}}
\toprule
\textbf{Método} & \textbf{Cuándo se usa y qué evidencia produce} \\
\midrule
\textbf{Prueba} & Prueba automatizada (unitaria o de integración) que se ejecuta en CI y falla
si el criterio deja de cumplirse. Evidencia: la clase de prueba, declarada en la matriz. \\
\textbf{Demostración} & Ejecución observada del sistema desplegado sobre la topología real.
Evidencia: la transcripción o captura archivada en \cd{docs/mediciones/} o
\cd{docs/evidencias/}, con fecha. \\
\textbf{Análisis} & Medición instrumentada o cálculo sobre datos crudos archivados
(k6, JaCoCo, Lighthouse, SUS, ZAP). Evidencia: el archivo crudo y el comando que lo generó,
mapeados en \cd{docs/mediciones/DATA-PROVENANCE.md}. \\
\textbf{Inspección} & Revisión de código, configuración o migración cuando el criterio es
estructural y no se puede ejercitar en tiempo de ejecución. Evidencia: el artefacto inspeccionado
y la fecha de la inspección. \\
\bottomrule
\end{tabular}
\end{table}

\subsection{Referencias}
\begin{itemize}
    \item ISO/IEC/IEEE 29148:2018 --- \emph{Systems and software engineering --- Life cycle
          processes --- Requirements engineering}.
    \item INCOSE \emph{Guide for Writing Requirements}, v4.
    \item OWASP \emph{Top 10} (2021) --- referencia de los requisitos de seguridad.
    \item WCAG 2.1 nivel AA --- referencia del requisito de accesibilidad.
    \item Ley Orgánica de Protección de Datos Personales del Ecuador (LOPDP).
    \item Brooke, J. (1996). \emph{SUS: A quick and dirty usability scale} --- instrumento de
          RNF-22.
    \item Decisiones de arquitectura del proyecto: \cd{docs/adr/ADR-001} a \cd{ADR-007}.
\end{itemize}

% ======================= 3. DESCRIPCIÓN GENERAL ==============================
\section{Descripción general}

\subsection{Perspectiva del producto}
Solución web de arquitectura en N capas: cliente Angular (SPA) servido por nginx como proxy
inverso, servicio de aplicación Spring Boot que expone una API REST, base de datos PostgreSQL
con esquema \texttt{presus} versionado por migraciones Flyway y una parte de la lógica de
negocio resuelta en procedimientos almacenados PL/pgSQL, y Redis como caché, almacén de
\emph{refresh tokens}, lista de revocación de tokens de acceso y contador del límite de tasa.
Todos los componentes se orquestan con Docker Compose.

\textbf{Tamaño real del sistema en el punto de corte:} 31 controladores REST, 207 combinaciones
de método y ruta, 10 rutinas SQL (8 procedimientos y 2 funciones de auditoría), 30 migraciones
Flyway y 559 pruebas automatizadas. Estas cifras están contadas sobre la etiqueta
\texttt{v1.0.1} y se recuentan en cada versión de este documento; no se dejan sin fecha.

\subsection{Actores del sistema}
\label{sec:actores}

\begin{table}[H]\centering\small
\begin{tabular}{p{2.6cm} p{4.2cm} p{7.6cm}}
\toprule
\textbf{Actor} & \textbf{Descripción} & \textbf{Responsabilidad en el proceso} \\
\midrule
Estudiante & Postulante a titulación. & Registra su solicitud, entrega correcciones en las
tres fases de tutoría, carga el anteproyecto, consulta su ruta de progreso y su acta. \\
Docente & Personal académico. & Puede actuar como tutor, como miembro del tribunal, o ambos en
solicitudes distintas. Nunca ambos en la misma solicitud. \\
Tutor & Rol funcional de un Docente sobre una solicitud. & Abre y aprueba las fases de tutoría,
revisa el anteproyecto y firma el acta con rol \texttt{TUTOR}. \\
Miembro del tribunal & Rol funcional de un Docente sobre una solicitud
(\texttt{PRESIDENTE}, \cd{VOCAL\_1}, \cd{VOCAL\_2}). & Registra su evaluación por
criterio de rúbrica y firma el acta con su rol. \\
Coordinador & Gestor académico de la carrera. & Revisa y aprueba solicitudes, asigna tutor y
tribunal, programa el cronograma, gestiona salas y rúbricas, y consulta reportes. \\
Administrador & Administrador del sistema. & Gestiona usuarios, roles y permisos, catálogos
académicos, la bitácora de auditoría y los respaldos. \\
\bottomrule
\end{tabular}
\caption{Actores. \emph{Tutor} y \emph{miembro del tribunal} no son roles de cuenta sino roles
funcionales asignados por solicitud: un mismo Docente los ejerce o no según la solicitud. Esta
distinción es la que explica por qué un permiso no basta para firmar un acta (ver RF-37).}
\end{table}

\subsection{Restricciones de diseño}
\label{sec:restricciones}

Estas son imposiciones sobre la solución, no requisitos del sistema. Se declaran aparte porque
la norma distingue requisito de restricción, y porque enunciarlas como requisitos («el sistema
usa PostgreSQL») describiría lo construido en vez de especificar lo exigido.

\begin{table}[H]\centering\small
\begin{tabular}{p{1.5cm} p{13.2cm}}
\toprule
\textbf{ID} & \textbf{Restricción y origen} \\
\midrule
RD-01 & La persistencia debe ser PostgreSQL 15 o superior, con el esquema versionado por
migraciones Flyway y \texttt{ddl-auto=validate}. Origen: ADR-006 (estrategia híbrida
ORM + procedimientos almacenados). \\
RD-02 & El cliente debe ser una SPA Angular servida por nginx; el navegador nunca consume
servicios de terceros directamente. Origen: ADR-003. \\
RD-03 & La autenticación debe ser sin estado de sesión en el servidor, basada en JWT.
Origen: ADR-002. \\
RD-04 & El despliegue debe ser reproducible con Docker Compose desde una clonación limpia.
Origen: ADR-007 y criterio de reproducibilidad de la asignatura. \\
RD-05 & La autorización debe resolverse contra un catálogo de permisos almacenado en base de
datos, no contra roles fijos en código. Origen: ADR-005. \\
\bottomrule
\end{tabular}
\end{table}

\subsection{Supuestos y dependencias}

\begin{enumerate}
    \item Las cuentas las provisiona el Administrador. \textbf{No existe autorregistro público},
          y esto es una decisión deliberada, no una omisión: el sistema trata expedientes
          académicos nominales.
    \item El correo institucional es el identificador único de la cuenta.
    \item El sistema depende de PostgreSQL y de Redis. El comportamiento exigido ante la caída
          de Redis se especifica en RNF-04 y \textbf{no} se deja como nota de riesgo: una
          política de degradación es una decisión de seguridad y merece un requisito.
    \item El único servicio externo consumido es el directorio de universidades descrito en la
          sección~\ref{sec:iface-consumidas}. \textbf{No se envía ningún dato personal a
          terceros}, y RNF-20 lo convierte en obligación exigible para que ninguna refactorización
          futura pueda cambiarlo sin romper un requisito.
    \item El envío de correo saliente está previsto pero desactivado por configuración
          (\cd{app.mail.enabled}). Mientras esté desactivado, las notificaciones son
          exclusivamente internas (RF-41).
\end{enumerate}

% ======================= 4. INTERFACES EXTERNAS ==============================
\section{Interfaces externas}
\label{sec:interfaces}

La norma distingue la interfaz que el sistema \emph{expone} de la que \emph{consume}. Se
separan aquí porque son contratos distintos: en la expuesta el sistema fija las reglas, y en la
consumida las sufre, de modo que cada interfaz consumida necesita además un comportamiento
declarado ante la indisponibilidad del otro extremo.

\subsection{Interfaces expuestas}

\begin{table}[H]\centering\small
\begin{tabular}{p{2.3cm} p{2.8cm} p{3.2cm} p{5.3cm}}
\toprule
\textbf{Interfaz} & \textbf{Protocolo / formato} & \textbf{Autenticación} & \textbf{Contrato} \\
\midrule
API REST & HTTP/1.1, JSON UTF-8, sobre estructura \texttt{ResponseWrapper} \{\texttt{success},
\texttt{message}, \texttt{data}\} & \emph{Bearer} JWT en cabecera \texttt{Authorization}, o
\emph{cookie} \texttt{jwtToken} \texttt{HttpOnly} & 207 combinaciones método+ruta bajo
\cd{/api/} y \cd{/api/v1/}; contrato publicado en OpenAPI~3.0 en \cd{/v3/api-docs} y
\cd{/swagger-ui/index.html}. \\
\addlinespace
Sonda de salud & HTTP, JSON & Ninguna (pública) & \cd{/actuator/health} con detalle por
componente (\texttt{db}, \texttt{redis}, \texttt{diskSpace}); es lo que permite comprobar que
las migraciones se aplicaron sin autenticarse. \\
\addlinespace
Proxy inverso & HTTP/HTTPS & N/A & nginx sirve la SPA y reenvía \cd{/api/} al servicio de
aplicación, con las mismas cabeceras de seguridad que exige RNF-10. \\
\bottomrule
\end{tabular}
\end{table}

\subsection{Interfaces consumidas}
\label{sec:iface-consumidas}

\begin{table}[H]\centering\small
\begin{tabular}{p{2.6cm} p{2.6cm} p{2.4cm} p{6.4cm}}
\toprule
\textbf{Proveedor} & \textbf{Protocolo} & \textbf{Auten\-ti\-ca\-ción} & \textbf{Comportamiento
exigido ante indisponibilidad} \\
\midrule
Directorio de universidades (\emph{Hipo Labs}) & HTTP, JSON, \cd{GET /search?country=Ecuador} &
Ninguna & Tiempo de espera de 5\,s en conexión, lectura y escritura; 3 reintentos con retroceso
exponencial desde 1\,s; agotados, respuesta de reserva local. Nunca propaga el fallo al cliente
(RF-45). \\
\addlinespace
PostgreSQL & TCP, JDBC & Usuario y contraseña por variable de entorno & Indisponibilidad
total del servicio. La aplicación no debe arrancar si las migraciones no se aplican. \\
\addlinespace
Redis & TCP (RESP) & Contraseña opcional por variable de entorno & Degradación explícita y
\textbf{fail-closed} en la verificación de revocación de tokens y en el límite de tasa
(RNF-04). \\
\addlinespace
SMTP saliente & SMTP/TLS & Usuario y contraseña por variable de entorno & Desactivado por
defecto (\cd{app.mail.enabled=false}). Con el envío activo, un fallo del proveedor no debe
impedir que la notificación interna se registre (RF-41). \\
\bottomrule
\end{tabular}
\end{table}

\noindent\textbf{Nota deliberada sobre proveedores de inteligencia artificial:} el asistente
virtual (RF-44) y la generación de ideas de tema (RF-47) se resuelven \textbf{íntegramente
dentro del sistema}, sin consumir ningún proveedor externo de modelo de lenguaje. Ningún dato
personal sale hacia un tercero. RNF-20 convierte esa propiedad en una obligación verificable,
porque hoy vive solo en el código y nada impediría que una refactorización futura empezara a
enviar datos.

% ===================== 5. ESTADOS DEL DOMINIO ================================
\section{Catálogo de estados del dominio y transiciones}
\label{sec:estados}

Los criterios de aceptación de la sección~\ref{sec:rf} usan nombres de estado. Esta sección
declara, para cada entidad con ciclo de vida, el \textbf{conjunto cerrado} de sus valores
admisibles y las transiciones válidas con el actor o evento que las dispara. Sin esta tabla los
criterios de aceptación se apoyan en vocabulario no definido.

\subsection{Solicitud de pre-sustentación}

\begin{table}[H]\centering\small
\begin{tabular}{p{2.4cm} p{4.6cm} p{7.4cm}}
\toprule
\textbf{Estado} & \textbf{Significado} & \textbf{Transición de salida (disparador)} \\
\midrule
\texttt{CREADA} & Registrada, aún no enviada. & $\rightarrow$ \texttt{ENVIADA} (Estudiante
envía). No es suspendible. \\
\texttt{ENVIADA} & En revisión de la coordinación. & $\rightarrow$ \texttt{APROBADA} /
\texttt{RECHAZADA} (Coordinador). \\
\texttt{APROBADA} & Aceptada; habilita la asignación de tutor. & $\rightarrow$ \texttt{TUTORIA}
(asignación de tutor). \\
\texttt{TUTORIA} & Acompañamiento en curso. & $\rightarrow$ \texttt{EVALUACION} (asignación de
tribunal, que exige tutoría completada). \\
\texttt{EVALUACION} & Tribunal conformado y defensa programable. & $\rightarrow$
\texttt{CALIFICADA} (registro de la nota final). \\
\texttt{CALIFICADA} & Nota final registrada; habilita el acta. & $\rightarrow$
\texttt{COMPLETADA} (última firma del acta). \\
\texttt{COMPLETADA} & Proceso cerrado. & Terminal. \\
\texttt{RECHAZADA} & Rechazada por la coordinación. & Terminal. No es suspendible. \\
\texttt{SUSPENDIDA} & Detenida por decisión académica, con motivo obligatorio. & Terminal en
esta versión. No es suspendible de nuevo. \\
\bottomrule
\end{tabular}
\caption{Estados de la solicitud. Cualquier estado distinto de \texttt{CREADA},
\texttt{RECHAZADA} y \texttt{SUSPENDIDA} admite la transición a \texttt{SUSPENDIDA} (RF-16).}
\end{table}

\subsection{Acta de pre-sustentación}

\begin{table}[H]\centering\small
\begin{tabular}{p{2.4cm} p{5.0cm} p{7.0cm}}
\toprule
\textbf{Estado} & \textbf{Significado} & \textbf{Transiciones válidas de salida} \\
\midrule
\texttt{GENERADA} & Acta emitida en PDF. & \texttt{REVISADA}, \texttt{OBSERVADA},
\texttt{ANULADA}. \\
\texttt{REVISADA} & Revisada académicamente. & \texttt{FINALIZADA}, \texttt{OBSERVADA},
\texttt{ANULADA}. \\
\texttt{OBSERVADA} & Con observaciones; exige motivo. & \texttt{REVISADA}, \texttt{GENERADA},
\texttt{ANULADA}. \\
\texttt{FINALIZADA} & Cerrada. Se alcanza también de forma automática con la cuarta firma. &
\texttt{ANULADA}. \\
\texttt{ANULADA} & Anulada; exige motivo. & Terminal. \\
\bottomrule
\end{tabular}
\caption{Estados del acta. \texttt{OBSERVADA} y \texttt{ANULADA} exigen motivo. La transición a
\texttt{ANULADA} desde cualquier estado está reservada al Administrador (RF-39).}
\end{table}

\subsection{Tutoría, fase de tutoría y demás entidades}

\begin{table}[H]\centering\small
\begin{tabular}{p{3.1cm} p{5.4cm} p{6.0cm}}
\toprule
\textbf{Entidad} & \textbf{Conjunto cerrado de estados} & \textbf{Disparador del cambio} \\
\midrule
Tutoría (\texttt{tutores}) & \texttt{ACTIVO}, \texttt{COMPLETADA} & \texttt{COMPLETADA} cuando
las 3 fases quedan aprobadas (RF-26). \\
Fase de tutoría & \cd{PENDIENTE\_TUTOR}, \cd{PENDIENTE\_ESTUDIANTE},
\texttt{APROBADA} & El tutor abre y aprueba; el estudiante entrega. \\
Mensaje de fase & \texttt{OBSERVACION}, \texttt{RESPUESTA}, \texttt{APROBACION} & Tipo fijado
por el emisor y la acción. \\
Anteproyecto (\cd{estados\_proceso}) & \texttt{PENDIENTE}, \cd{EN\_PROCESO},
\texttt{APROBADO}, \texttt{OBSERVADO}, \texttt{RECHAZADO} & Revisión del tutor (RF-21). \\
Cronograma & \texttt{PROGRAMADO} & Único estado en uso en esta versión. \\
Resultado de evaluación & \texttt{APROBADO} ($\geq 7{,}0$), \texttt{REPROBADO} ($< 7{,}0$) &
Calculado, no editable (RF-34). \\
Estudiante (académico) & \texttt{ACTIVO}, \texttt{EGRESADO}, \texttt{GRADUADO},
\texttt{RETIRADO}, \texttt{SUSPENDIDO} & Gestión de estudiantes (RF-54). \\
Cuenta de usuario & \texttt{activo} verdadero / falso & Activación y desactivación (RF-09). \\
\bottomrule
\end{tabular}
\end{table}

\noindent\fbox{\parbox{0.98\textwidth}{\textbf{Hallazgo declarado (2026-09-09).} Los nueve
estados de la solicitud y el único estado del cronograma \textbf{no están sembrados por ninguna
migración}: se crean bajo demanda la primera vez que el proceso los necesita. La consecuencia es
que en una base recién migrada el catálogo está vacío, y tanto el identificador como el campo
\texttt{orden} de cada estado dependen del orden en que se ejecuten los casos de uso. RF-13
declara el conjunto cerrado y RNF-25 exige que la siembra sea determinista; hasta que esa
siembra exista, el catálogo de estados no es reproducible entre dos despliegues. Se declara aquí
en vez de presentar la tabla como si fuera un catálogo sembrado.}}

% ================== 6. MATRIZ DE PERMISOS ROL x OPERACIÓN ====================
\section{Matriz de permisos por rol y operación}
\label{sec:permisos}

Con 31 controladores y 207 rutas, los permisos ya no se pueden auditar leyendo requisito por
requisito: hace falta una tabla normativa. Esta sección la fija. Es la única forma de decidir si
una asimetría de acceso es una decisión de negocio o un defecto, y es lo que permite detectar un
permiso que se otorgó sin querer.

\subsection{Modelo de autorización}

El sistema define \textbf{4 roles de cuenta} (\texttt{ADMIN}, \texttt{COORDINADOR},
\texttt{DOCENTE}, \texttt{ESTUDIANTE}) y \textbf{29 permisos} en catálogo de base de datos. La
autorización se resuelve por permiso, no por rol, y la asignación permiso\,$\rightarrow$\,rol es
editable en caliente (RF-12). Sobre ese control se aplican, cuando corresponde, dos
comprobaciones adicionales que el permiso por sí solo no cubre:

\begin{itemize}
    \item \textbf{Propiedad del recurso:} el titular solo accede a lo suyo (RNF-13).
    \item \textbf{Rol funcional en la solicitud:} tener \cd{ACTA\_FIRMAR} no habilita a
          firmar cualquier acta; hay que ser además el tutor o el miembro del tribunal de esa
          solicitud (RF-37).
\end{itemize}

\subsection{Asignación normativa de permisos}

\begin{longtable}{p{5.1cm} p{1.6cm} p{1.9cm} p{2.2cm} p{1.7cm}}
\toprule
\textbf{Permiso} & \textbf{ADMIN} & \textbf{COORD.} & \textbf{DOCENTE} & \textbf{ESTUD.} \\
\midrule
\endfirsthead
\toprule
\textbf{Permiso} & \textbf{ADMIN} & \textbf{COORD.} & \textbf{DOCENTE} & \textbf{ESTUD.} \\
\midrule
\endhead
\multicolumn{5}{l}{\textit{Administración}}\\
\cd{USUARIOS\_GESTIONAR}            & Sí & --- & --- & --- \\
\cd{ROLES\_PERMISOS\_GESTIONAR}     & Sí & --- & --- & --- \\
\cd{AUDITORIA\_VER}                 & Sí & --- & --- & --- \\
\cd{BACKUPS\_GESTIONAR}             & Sí & --- & --- & --- \\
\cd{CARRERAS\_GESTIONAR}            & Sí & --- & --- & --- \\
\cd{ESTUDIANTES\_GESTIONAR}         & Sí & Sí  & --- & --- \\
\cd{NOTIFICACIONES\_GLOBAL\_VER}    & Sí & --- & --- & --- \\
\cd{NOTIFICACIONES\_ENVIAR}         & Sí & Sí  & --- & --- \\
\addlinespace
\multicolumn{5}{l}{\textit{Solicitudes y proceso}}\\
\cd{SOLICITUDES\_REVISAR}           & Sí & Sí  & Sí  & --- \\
\cd{SOLICITUDES\_SUSPENDER}         & Sí & Sí  & --- & --- \\
\cd{TRIBUNAL\_TUTOR\_ASIGNAR}       & Sí & Sí  & --- & --- \\
\cd{CRONOGRAMA\_GESTIONAR}          & Sí & Sí  & --- & --- \\
\cd{SALA\_GESTIONAR}                & Sí & Sí  & --- & --- \\
\addlinespace
\multicolumn{5}{l}{\textit{Tutoría y anteproyecto}}\\
\cd{ANTEPROYECTO\_REVISAR}          & Sí & Sí  & Sí  & --- \\
\cd{TUTORIA\_GESTIONAR}             & Sí & Sí  & Sí  & --- \\
\cd{TUTORIA\_AVANCE\_ESTUDIANTE}    & Sí & Sí  & --- & Sí  \\
\addlinespace
\multicolumn{5}{l}{\textit{Evaluación}}\\
\cd{RUBRICA\_GESTIONAR}             & Sí & Sí  & --- & --- \\
\cd{EVALUACION\_RUBRICA\_REGISTRAR} & Sí & Sí  & Sí  & --- \\
\cd{EVALUACION\_CALIFICAR}          & Sí & Sí  & --- & --- \\
\addlinespace
\multicolumn{5}{l}{\textit{Actas}}\\
\cd{ACTA\_GENERAR}                  & Sí & Sí  & --- & --- \\
\cd{ACTA\_FIRMAR}                   & Sí & Sí  & Sí  & --- \\
\cd{ACTAS\_VER}                     & Sí & Sí  & --- & --- \\
\cd{ACTAS\_VER\_PROPIAS}            & Sí & --- & Sí  & --- \\
\cd{ACTA\_HISTORIAL\_VER}           & Sí & Sí  & Sí  & --- \\
\cd{ACTA\_ESTADO\_CAMBIAR}          & Sí & Sí  & --- & --- \\
\cd{ACTAS\_GESTIONAR}               & Sí & --- & --- & --- \\
\addlinespace
\multicolumn{5}{l}{\textit{Reportes y orientación}}\\
\cd{REPORTES\_VER}                  & Sí & Sí  & --- & --- \\
\cd{ORIENTACION\_TEMAS\_VER}        & Sí & Sí  & Sí  & Sí  \\
\cd{ORIENTACION\_CATALOGO\_GESTIONAR} & Sí & Sí & --- & --- \\
\bottomrule
\caption{Asignación normativa permiso\,$\times$\,rol. «Sí» significa que el rol posee el permiso
en la configuración inicial sembrada; «---» que no lo posee. La asignación es editable en
caliente por RF-12, de modo que esta tabla especifica el \textbf{estado inicial exigido}, no una
constante inmutable.}
\end{longtable}

\subsection{Acceso que no depende del catálogo de permisos}

Tres capacidades se autorizan por rol de cuenta o por simple autenticación, no por permiso. Se
declaran aparte para que la tabla anterior no se lea como la superficie completa de
autorización:

\begin{table}[H]\centering\small
\begin{tabular}{p{4.4cm} p{3.4cm} p{6.8cm}}
\toprule
\textbf{Capacidad} & \textbf{Control} & \textbf{Requisito} \\
\midrule
Ruta de progreso de titulación & Rol \texttt{ESTUDIANTE} & RF-51 \\
Temas guardados & Rol \texttt{ESTUDIANTE} & RF-48 \\
Consulta de catálogos de apoyo, docentes, asistente virtual, universidades, estado en tiempo
real, notificaciones propias, perfil propio & Usuario autenticado & RF-53, RF-55, RF-44, RF-45,
RF-43, RF-41, RF-10 \\
\bottomrule
\end{tabular}
\end{table}

\subsection{Asimetrías declaradas}
\label{sec:asimetrias}

Tres asimetrías de la tabla anterior son intencionales y se explican aquí para que una revisión
futura no las tome por defectos, y una cuarta es una desviación real que queda como deuda:

\begin{enumerate}
    \item \textbf{El Coordinador no gestiona catálogos académicos}
          (\cd{CARRERAS\_GESTIONAR} es exclusivo del Administrador) aunque sí gestiona
          estudiantes. Es deliberado: facultades, carreras, modalidades y períodos son datos
          maestros institucionales.
    \item \textbf{El Coordinador no accede a la bitácora de auditoría}
          (\cd{AUDITORIA\_VER}) ni a los respaldos. La auditoría debe poder registrar
          acciones del Coordinador sin que él pueda consultarlas.
    \item \textbf{El Administrador posee \cd{ACTAS\_VER\_PROPIAS}} pese a no ser docente. Es
          consecuencia de otorgarle el conjunto completo; carece de efecto porque el filtro por
          propiedad no le devuelve actas.
    \item \textbf{Deuda declarada:} el Coordinador posee \cd{TUTORIA\_AVANCE\_ESTUDIANTE},
          \cd{ACTA\_FIRMAR}, \cd{ANTEPROYECTO\_REVISAR} y \cd{TUTORIA\_GESTIONAR}
          por herencia de la regla «todos los permisos menos tres» con la que se sembró el
          catálogo inicial, no por una decisión sobre cada uno. No hay riesgo de acceso indebido
          porque las comprobaciones de rol funcional de RF-37, RF-21 y RF-23 rechazan a quien no
          es el tutor o el jurado de esa solicitud; pero la asignación no refleja una decisión
          y debe revisarse permiso por permiso.
\end{enumerate}

% ==================== 7. REQUISITOS FUNCIONALES ==============================
\section{Requisitos funcionales}
\label{sec:rf}

Cada requisito se especifica con la plantilla completa que exige la
sección~\ref{sec:convencion}: enunciado normativo, prioridad MoSCoW, fuente, endpoint,
\emph{rationale}, criterio de aceptación observable, método de verificación
(sección~\ref{sec:metodos}) y estado del vocabulario cerrado (sección~\ref{sec:estados-req}).
Los endpoints se citan tal como los expone el servicio en la etiqueta \texttt{v1.0.1}; la ruta
de archivo y la clase de prueba viven en la matriz de trazabilidad, no aquí.

\subsection{Autenticación y gestión de la sesión}

\REQ{RF-01}{Autenticación con credenciales institucionales}
\campo{Enunciado}{El sistema debe autenticar a un usuario a partir de su correo institucional y
su contraseña, y emitir un token de acceso y un token de renovación asociados a esa sesión.}
\campo{Prioridad}{Must}
\campo{Fuente}{HU-01, CU-01, ADR-002}
\campo{Endpoint}{\cd{POST /api/v1/auth/login}}
\campo{Rationale}{Es la puerta de entrada de todo el sistema: ningún otro requisito es
alcanzable sin ella, y su diseño determina qué puede hacer un atacante que capture una
credencial.}
\campo{Criterio de aceptación}{}
\begin{itemize}
    \item Con credenciales válidas responde \texttt{200} y devuelve el token de acceso, el
          identificador, el correo, el nombre y el rol del usuario.
    \item El token de acceso lleva los siete atributos \texttt{jti}, \texttt{iss}, \texttt{sub},
          \texttt{aud}, \texttt{iat}, \texttt{nbf} y \texttt{exp}.
    \item Con contraseña incorrecta responde \texttt{401}. Con un correo no registrado responde
          \textbf{también \texttt{401}}, con el mismo cuerpo: la respuesta no debe permitir
          distinguir un correo existente de uno inexistente.
    \item Nunca devuelve la contraseña ni su representación almacenada en el cuerpo de la
          respuesta.
    \item Superado el límite de RNF-09 responde \texttt{429}.
\end{itemize}
\campo{Método de verificación}{Prueba}
\campo{Estado}{\estadoV}
\campo{Depende de}{RNF-05, RNF-07, RNF-09}

\REQ{RF-02}{Renovación del token de acceso con rotación del token de renovación}
\campo{Enunciado}{El sistema debe emitir un nuevo token de acceso a partir de un token de
renovación válido, invalidando en el mismo acto el token de renovación empleado y emitiendo uno
nuevo.}
\campo{Prioridad}{Must}
\campo{Fuente}{ADR-002; OWASP A07 (fallos de identificación y autenticación)}
\campo{Endpoint}{\cd{POST /api/v1/auth/refresh}}
\campo{Rationale}{Sin rotación, un token de renovación robado sirve indefinidamente durante toda
su vida útil. Con rotación, su reutilización se vuelve detectable, que es lo que convierte el
robo en un evento observable en vez de silencioso.}
\campo{Criterio de aceptación}{}
\begin{itemize}
    \item Con un token de renovación vigente responde \texttt{200} con un token de acceso nuevo
          y un token de renovación nuevo.
    \item El token de renovación empleado deja de ser válido de inmediato: un segundo uso
          responde \texttt{401}.
    \item Si el token presentado es uno ya utilizado, el sistema debe tratarlo como
          \textbf{reutilización} y responder \texttt{401} \textbf{revocando además todas las
          sesiones activas de ese usuario}, no solo rechazando la petición.
    \item Sin token de renovación en la petición responde \texttt{400}.
\end{itemize}
\campo{Método de verificación}{Prueba}
\campo{Estado}{\estadoI}
\campo{Nota}{La detección de reutilización está construida y es observable, pero no hay todavía
una prueba automatizada que ejercite el escenario de reutilización de punta a punta. Se declara
\estadoI{} y no \estadoV{} en vez de dar por verificado lo que la matriz no respalda.}

\REQ{RF-03}{Cierre de sesión con revocación del token de acceso}
\campo{Enunciado}{El sistema debe invalidar, al cerrar sesión, tanto el token de renovación de
esa sesión como el token de acceso presentado, de modo que este último deje de aceptarse aunque
no haya expirado.}
\campo{Prioridad}{Must}
\campo{Fuente}{ADR-002}
\campo{Endpoint}{\cd{POST /api/v1/auth/logout}}
\campo{Rationale}{Un token de acceso sin estado es válido hasta su expiración por construcción.
Sin una lista de revocación consultada en cada petición, «cerrar sesión» sería una acción
puramente cosmética en el cliente.}
\campo{Criterio de aceptación}{}
\begin{itemize}
    \item Responde \texttt{200} y limpia las \emph{cookies} de sesión del cliente.
    \item Una petición posterior a cualquier endpoint protegido con \textbf{ese mismo} token de
          acceso responde \texttt{401}, aun estando dentro de su ventana de validez.
    \item La entrada de revocación debe expirar por sí sola al vencer el token, sin crecimiento
          ilimitado del almacén.
\end{itemize}
\campo{Método de verificación}{Prueba}
\campo{Estado}{\estadoV}
\campo{Depende de}{RNF-04 (la revocación se apoya en Redis y su política de degradación es la
que decide qué pasa si el almacén no responde)}

\REQ{RF-04}{Provisión de cuentas por el Administrador}
\campo{Enunciado}{El sistema debe permitir únicamente al Administrador crear cuentas de usuario,
aceptando exclusivamente nombre, apellido, correo, contraseña y rol.}
\campo{Prioridad}{Must}
\campo{Fuente}{HU-11, CU-11; hallazgo de auditoría OWASP A01 del 2026-09-04}
\campo{Endpoints}{\cd{POST /api/v1/auth/register}, \cd{POST /api/usuarios}}
\campo{Rationale}{El sistema trata expedientes académicos nominales, así que el alta no puede
ser pública. Además, aceptar la entidad completa del cliente permitiría fijar por asignación
masiva campos que no le corresponden (identificador, estado activo, rol interno, fecha de
creación).}
\campo{Criterio de aceptación}{}
\begin{itemize}
    \item Sin el permiso \cd{USUARIOS\_GESTIONAR} responde \texttt{403}.
    \item Ningún campo del cuerpo distinto de los cinco admitidos debe alterar el registro
          creado, aunque se envíe.
    \item Con un correo ya registrado responde \texttt{400} sin crear el usuario.
    \item El rol declarado y el rol interno del usuario quedan siempre resueltos al mismo valor;
          no es posible crear una cuenta con ambos desincronizados.
    \item La contraseña se almacena conforme a RNF-05 y nunca se devuelve.
\end{itemize}
\campo{Método de verificación}{Prueba}
\campo{Estado}{\estadoV}
\campo{Depende de}{RNF-06}

\REQ{RF-05}{Recuperación de la contraseña olvidada}
\campo{Enunciado}{El sistema debe permitir al titular de una cuenta solicitar la recuperación de
su contraseña indicando su correo institucional, y restablecerla mediante un elemento de un solo
uso con vigencia limitada enviado a ese correo.}
\campo{Prioridad}{Must}
\campo{Fuente}{Auditoría de cobertura del 2026-09-09; el asistente virtual (RF-44) ya ofrece
esta capacidad al usuario. Historia de usuario pendiente (\S\ref{sec:pendientes}).}
\campo{Endpoints}{Por definir. No existen en el punto de corte.}
\campo{Rationale}{Es la única carencia del corpus que exige programar y no solo redactar. Como
las cuentas las crea exclusivamente el Administrador (RF-04) y la actualización de usuario
\textbf{no modifica la contraseña}, hoy un usuario que la olvida no tiene ninguna vía de
recuperación dentro del sistema: ni propia, ni administrativa. La única salida es una
intervención directa sobre la base de datos, que no está especificada, no deja traza de
autorización y no es auditable.}
\campo{Criterio de aceptación}{}
\begin{itemize}
    \item La solicitud responde \texttt{200} con el mismo cuerpo tanto si el correo está
          registrado como si no: la respuesta no debe permitir enumerar cuentas.
    \item El elemento de un solo uso caduca a los 30 minutos y no puede emplearse dos veces.
    \item Al restablecerse la contraseña se aplican las exigencias de RNF-06 y \textbf{se
          revocan todas las sesiones activas} del titular.
    \item El número de solicitudes por cuenta y por ventana de tiempo está limitado, para que el
          endpoint no sea un vector de abuso de correo.
\end{itemize}
\campo{Método de verificación}{Prueba (cuando exista la funcionalidad)}
\campo{Estado}{\estadoP}
\campo{Nota}{Se especifica antes de construirse, y a propósito: sin requisito no hay criterio
contra el cual construirlo. Alternativa admisible si el equipo decide no implementarlo en esta
entrega: declararlo explícitamente fuera de alcance y especificar en su lugar el restablecimiento
administrativo, que hoy tampoco existe.}

\REQ{RF-06}{Cambio de la contraseña propia}
\campo{Enunciado}{El sistema debe permitir a un usuario autenticado cambiar su propia contraseña
presentando la contraseña vigente.}
\campo{Prioridad}{Must}
\campo{Fuente}{Auditoría de cobertura del 2026-09-09. Historia de usuario pendiente.}
\campo{Endpoint}{Por definir. No existe en el punto de corte.}
\campo{Rationale}{Las contraseñas iniciales las fija el Administrador (RF-04), de modo que sin
esta capacidad el usuario nunca controla su propia credencial y una contraseña inicial conocida
por un tercero permanece válida indefinidamente.}
\campo{Criterio de aceptación}{}
\begin{itemize}
    \item Con la contraseña vigente correcta responde \texttt{200}; con la incorrecta,
          \texttt{401}, sin cambiar nada.
    \item La contraseña nueva cumple RNF-06 y no puede ser igual a la vigente.
    \item El cambio revoca todas las sesiones activas del titular salvo aquella desde la que se
          realiza.
\end{itemize}
\campo{Método de verificación}{Prueba (cuando exista la funcionalidad)}
\campo{Estado}{\estadoP}

\REQ{RF-07}{Consulta de los permisos efectivos propios}
\campo{Enunciado}{El sistema debe devolver al usuario autenticado el conjunto de códigos de
permiso que su rol tiene concedidos en ese momento.}
\campo{Prioridad}{Should}
\campo{Fuente}{ADR-005}
\campo{Endpoint}{\cd{GET /api/me/permisos}}
\campo{Rationale}{Permite al cliente ocultar lo que el usuario no puede ejecutar sin duplicar el
mapa de permisos en el frontend. La decisión de autorización sigue siendo del servidor
(RNF-12): esto es presentación, no control de acceso.}
\campo{Criterio de aceptación}{}
\begin{itemize}
    \item Responde \texttt{200} con la lista de códigos de permiso del usuario en sesión.
    \item El resultado refleja los cambios de RF-12 sin necesidad de volver a iniciar sesión.
    \item Sin autenticación responde \texttt{401}.
    \item Devuelve exclusivamente los permisos del usuario en sesión, nunca los de otro.
\end{itemize}
\campo{Método de verificación}{Prueba}
\campo{Estado}{\estadoI}

\subsection{Identidades y autorización}

\REQ{RF-08}{Gestión de usuarios}
\campo{Enunciado}{El sistema debe permitir al Administrador consultar, crear, modificar y
eliminar usuarios, y consultarlos de forma paginada, por correo y filtrados por estado activo.}
\campo{Prioridad}{Must}
\campo{Fuente}{HU-11, CU-11}
\campo{Endpoints}{\cd{GET /api/usuarios}, \cd{GET /api/usuarios/paginado},
\cd{GET /api/usuarios/\{id\}}, \cd{GET /api/usuarios/email/\{email\}},
\cd{GET /api/usuarios/activos}, \cd{PUT /api/usuarios/\{id\}}, \cd{DELETE /api/usuarios/\{id\}}}
\campo{Rationale}{Es el requisito que sostiene la identidad de todos los demás: sin gestión de
usuarios no hay a quién asignar tutor, tribunal ni acta.}
\campo{Criterio de aceptación}{}
\begin{itemize}
    \item Toda operación de esta capacidad sin el permiso \cd{USUARIOS\_GESTIONAR} responde
          \texttt{403}.
    \item Ninguna respuesta incluye la contraseña almacenada, en ninguna de las operaciones.
    \item El borrado de un usuario con historial académico asociado se rechaza con \texttt{400}
          en vez de dejar registros huérfanos; la vía prevista para ese caso es RF-09.
    \item La modificación de un usuario \textbf{no altera su contraseña}.
\end{itemize}
\campo{Método de verificación}{Prueba}
\campo{Estado}{\estadoV}
\campo{Nota}{El último criterio documenta el comportamiento real y es también el que hace
inevitable RF-05: hoy nadie, ni siquiera el Administrador, puede restablecer una contraseña
desde el sistema.}

\REQ{RF-09}{Activación y desactivación de cuentas}
\campo{Enunciado}{El sistema debe permitir al Administrador desactivar y reactivar una cuenta
sin eliminarla, y una cuenta desactivada no debe poder autenticarse.}
\campo{Prioridad}{Should}
\campo{Fuente}{HU-11}
\campo{Endpoints}{\cd{PATCH /api/usuarios/\{id\}/activar},
\cd{PATCH /api/usuarios/\{id\}/desactivar}}
\campo{Rationale}{Es la baja correcta cuando el usuario ya tiene historial académico: preserva
la trazabilidad de sus actas y evaluaciones, cosa que el borrado destruiría.}
\campo{Criterio de aceptación}{}
\begin{itemize}
    \item Responde \texttt{200} y el estado activo del usuario queda en el valor solicitado.
    \item Un intento de autenticación de una cuenta desactivada responde \texttt{401}.
    \item Sin el permiso \cd{USUARIOS\_GESTIONAR} responde \texttt{403}.
    \item El cambio queda registrado en la bitácora de auditoría con el actor que lo realizó
          (RNF-18).
\end{itemize}
\campo{Método de verificación}{Prueba}
\campo{Estado}{\estadoV}

\REQ{RF-10}{Actualización del perfil propio}
\campo{Enunciado}{El sistema debe permitir a un usuario autenticado actualizar su correo de
notificaciones y su teléfono, y solo los suyos.}
\campo{Prioridad}{Should}
\campo{Fuente}{Auditoría de cobertura del 2026-09-09. Historia de usuario pendiente.}
\campo{Endpoint}{\cd{PATCH /api/usuarios/\{id\}/perfil}}
\campo{Rationale}{El correo institucional identifica la cuenta y no debe poder cambiarlo el
propio usuario; el correo de notificaciones y el teléfono son datos de contacto y sí le
pertenecen.}
\campo{Criterio de aceptación}{}
\begin{itemize}
    \item Un intento de editar el perfil de otro usuario responde \texttt{403}, aunque quien lo
          intente sea Administrador.
    \item La operación no modifica el correo institucional, el rol ni el estado activo, aunque
          se envíen en el cuerpo.
\end{itemize}
\campo{Método de verificación}{Prueba}
\campo{Estado}{\estadoV}
\campo{Depende de}{RNF-13}

\REQ{RF-11}{Gestión de roles}
\campo{Enunciado}{El sistema debe permitir al Administrador consultar, crear, renombrar y
eliminar roles de usuario, conservando invariable el código del rol.}
\campo{Prioridad}{Should}
\campo{Fuente}{ADR-005}
\campo{Endpoints}{\cd{GET /api/roles}, \cd{POST /api/roles}, \cd{PUT /api/roles/\{id\}},
\cd{DELETE /api/roles/\{id\}}}
\campo{Rationale}{El nombre del rol es una etiqueta de presentación; su código es la clave que
usan las reglas de autorización. Permitir renombrar sin tocar el código deja adaptar el
vocabulario institucional sin romper el control de acceso.}
\campo{Criterio de aceptación}{}
\begin{itemize}
    \item Renombrar un rol devuelve \texttt{200} y el código del rol permanece inalterado.
    \item Eliminar un rol con usuarios asignados se rechaza con \texttt{400}.
    \item Sin el permiso \cd{ROLES\_PERMISOS\_GESTIONAR} responde \texttt{403}.
\end{itemize}
\campo{Método de verificación}{Prueba}
\campo{Estado}{\estadoV}

\REQ{RF-12}{Asignación de permisos a un rol}
\campo{Enunciado}{El sistema debe permitir al Administrador consultar el catálogo completo de
permisos y sustituir el conjunto de permisos concedidos a un rol, con efecto inmediato sobre las
sesiones activas.}
\campo{Prioridad}{Must}
\campo{Fuente}{ADR-005}
\campo{Endpoints}{\cd{GET /api/permisos}, \cd{PUT /api/permisos/rol/\{rolId\}}}
\campo{Rationale}{Es lo que hace que la matriz de la sección~\ref{sec:permisos} sea configuración
y no código: permite corregir una asignación indebida sin recompilar ni desplegar.}
\campo{Criterio de aceptación}{}
\begin{itemize}
    \item Responde \texttt{200} y el conjunto de permisos del rol queda exactamente en el
          enviado, sin residuos de la asignación anterior.
    \item Tras el cambio, una petición de un usuario de ese rol a un endpoint cuyo permiso se
          retiró responde \texttt{403} \textbf{sin necesidad de volver a iniciar sesión}.
    \item Sin el permiso \cd{ROLES\_PERMISOS\_GESTIONAR} responde \texttt{403}.
    \item El cambio queda registrado en la bitácora de auditoría (RNF-18).
\end{itemize}
\campo{Método de verificación}{Prueba}
\campo{Estado}{\estadoV}
\campo{Nota}{Este requisito puede dejar el sistema sin ningún usuario capaz de administrar
permisos si se retira \cd{ROLES\_PERMISOS\_GESTIONAR} al rol \texttt{ADMIN}. Esa
salvaguarda no está especificada ni construida, y queda como deuda declarada
(\S\ref{sec:pendientes}).}

\subsection{Solicitud de pre-sustentación}

\REQ{RF-13}{Registro de la solicitud de pre-sustentación}
\campo{Enunciado}{El sistema debe registrar la solicitud de pre-sustentación de un estudiante
con el título del tema, la modalidad de titulación y la convocatoria, dejándola en estado
\texttt{CREADA}.}
\campo{Prioridad}{Must}
\campo{Fuente}{HU-02, CU-02}
\campo{Endpoints}{\cd{POST /api/solicitudes/crear-por-usuario/\{usuarioId\}},
\cd{POST /api/solicitudes/crear/\{estudianteId\}}}
\campo{Rationale}{Es el hecho que abre el expediente: todo el resto del proceso cuelga de esta
solicitud. Registrarla como \texttt{CREADA} y no como \texttt{ENVIADA} separa la redacción del
compromiso, y deja al estudiante corregir antes de someterla.}
\campo{Criterio de aceptación}{}
\begin{itemize}
    \item El título del tema es obligatorio y no excede 300 caracteres; en otro caso responde
          \texttt{400} sin crear la solicitud.
    \item La modalidad de titulación es obligatoria; sin ella responde \texttt{400}.
    \item Si no se indica convocatoria, el sistema debe resolver la convocatoria activa; si no
          hay ninguna activa responde \texttt{400}.
    \item Si se indica un área temática que no pertenece a la línea de investigación declarada,
          responde \texttt{400}.
    \item La solicitud creada queda en estado \texttt{CREADA} conforme a la
          sección~\ref{sec:estados}, y al perfil de estudiante se le genera su código de
          expediente.
\end{itemize}
\campo{Método de verificación}{Prueba}
\campo{Estado}{\estadoV}

\REQ{RF-14}{Envío de la solicitud a revisión}
\campo{Enunciado}{El sistema debe permitir al estudiante propietario enviar su solicitud a
revisión, pasándola de \texttt{CREADA} a \texttt{ENVIADA}.}
\campo{Prioridad}{Must}
\campo{Fuente}{HU-02, CU-02}
\campo{Endpoint}{\cd{POST /api/solicitudes/enviar/\{id\}}}
\campo{Rationale}{Es el acto por el cual el estudiante asume la solicitud como definitiva y la
somete a la coordinación; sin él, la coordinación no puede distinguir un borrador de una
postulación.}
\campo{Criterio de aceptación}{}
\begin{itemize}
    \item Enviar una solicitud ajena responde \texttt{403}.
    \item La solicitud queda en estado \texttt{ENVIADA} y el cambio se registra en el historial
          de estados con el usuario que lo realizó.
    \item Enviar una solicitud que no esté en \texttt{CREADA} responde \texttt{400}.
\end{itemize}
\campo{Método de verificación}{Prueba}
\campo{Estado}{\estadoV}
\campo{Depende de}{RF-13}

\REQ{RF-15}{Aprobación o rechazo de la solicitud}
\campo{Enunciado}{El sistema debe permitir a quien posea el permiso
\cd{SOLICITUDES\_REVISAR} aprobar o rechazar una solicitud enviada, y rechazarla dejando una
observación dirigida al estudiante.}
\campo{Prioridad}{Must}
\campo{Fuente}{HU-02, CU-02}
\campo{Endpoints}{\cd{POST /api/solicitudes/aprobar/\{id\}},
\cd{POST /api/solicitudes/rechazar/\{id\}},
\cd{POST /api/solicitudes/rechazar-con-observacion/\{id\}}}
\campo{Rationale}{Un rechazo sin motivo obliga al estudiante a adivinar qué corregir. Separar el
rechazo con observación del rechazo simple permite exigir el motivo donde importa sin bloquear
el caso en que la solicitud es inadmisible de plano.}
\campo{Criterio de aceptación}{}
\begin{itemize}
    \item Aprobar deja la solicitud en \texttt{APROBADA}; rechazar, en \texttt{RECHAZADA}.
    \item El rechazo con observación exige el texto de la observación; sin él responde
          \texttt{400}.
    \item Sin el permiso \cd{SOLICITUDES\_REVISAR} responde \texttt{403}.
    \item Toda transición queda registrada en el historial de estados con usuario, estado
          anterior, estado nuevo y fecha.
    \item El estudiante recibe una notificación con el resultado (RF-41).
\end{itemize}
\campo{Método de verificación}{Prueba}
\campo{Estado}{\estadoV}

\REQ{RF-16}{Suspensión de una solicitud}
\campo{Enunciado}{El sistema debe permitir a quien posea el permiso
\cd{SOLICITUDES\_SUSPENDER} suspender una solicitud en curso indicando el motivo, y una
solicitud suspendida no debe admitir cargas de archivo por parte del estudiante.}
\campo{Prioridad}{Should}
\campo{Fuente}{Auditoría de cobertura del 2026-09-09. Historia de usuario pendiente.}
\campo{Endpoint}{\cd{POST /api/solicitudes/suspender/\{id\}}}
\campo{Rationale}{Detener un proceso sin motivo registrado deja al estudiante sin explicación y
a la institución sin traza de la decisión. La suspensión es reversible en su intención y
distinta del rechazo, que cierra el expediente.}
\campo{Criterio de aceptación}{}
\begin{itemize}
    \item El motivo de suspensión es obligatorio; sin él responde \texttt{400}.
    \item Solo son suspendibles las solicitudes en un estado distinto de \texttt{CREADA},
          \texttt{RECHAZADA} y \texttt{SUSPENDIDA}; en otro caso responde \texttt{400}.
    \item Tras la suspensión, la carga del anteproyecto (RF-19) responde \texttt{400} indicando
          el motivo registrado.
    \item Sin el permiso \cd{SOLICITUDES\_SUSPENDER} responde \texttt{403}.
\end{itemize}
\campo{Método de verificación}{Prueba}
\campo{Estado}{\estadoI}
\campo{Nota}{La reanudación de una solicitud suspendida no está especificada ni construida: hoy
\cd{SUSPENDIDA} es terminal. Queda como deuda declarada (\S\ref{sec:pendientes}).}

\REQ{RF-17}{Consulta y seguimiento de la solicitud propia}
\campo{Enunciado}{El sistema debe permitir al estudiante consultar sus propias solicitudes y el
seguimiento del avance de cada una, expresado como porcentaje y como etapa alcanzada.}
\campo{Prioridad}{Must}
\campo{Fuente}{HU-02, CU-02}
\campo{Endpoints}{\cd{GET /api/solicitudes/mis-solicitudes},
\cd{GET /api/solicitudes/\{id\}}, \cd{GET /api/solicitudes/\{id\}/seguimiento}}
\campo{Rationale}{El estudiante es el actor con menos visibilidad del proceso y el más afectado
por él; sin seguimiento propio, cada consulta de estado se convierte en un trámite presencial.}
\campo{Criterio de aceptación}{}
\begin{itemize}
    \item Devuelve exclusivamente las solicitudes del estudiante en sesión.
    \item Consultar el detalle o el seguimiento de una solicitud ajena responde \texttt{403}.
    \item El seguimiento devuelve la etapa alcanzada conforme a los estados de la
          sección~\ref{sec:estados} y un porcentaje coherente con esa etapa.
\end{itemize}
\campo{Método de verificación}{Prueba}
\campo{Estado}{\estadoV}
\campo{Depende de}{RNF-13}

\REQ{RF-18}{Consulta administrativa de solicitudes}
\campo{Enunciado}{El sistema debe permitir a quien posea el permiso
\cd{SOLICITUDES\_REVISAR} listar las solicitudes de forma completa y paginada, filtrarlas por
estudiante, contarlas por estado y obtener el reporte de defensas.}
\campo{Prioridad}{Should}
\campo{Fuente}{HU-09, CU-09}
\campo{Endpoints}{\cd{GET /api/solicitudes}, \cd{GET /api/solicitudes/paginado},
\cd{GET /api/solicitudes/contar-por-estado},
\cd{GET /api/solicitudes/estudiante/\{estudianteId\}},
\cd{GET /api/solicitudes/reporte-defensas}}
\campo{Rationale}{Es la vista de gestión de la coordinación. El conteo por estado alimenta los
indicadores de RF-56 sin recalcularlos en el cliente.}
\campo{Criterio de aceptación}{}
\begin{itemize}
    \item Sin el permiso \cd{SOLICITUDES\_REVISAR} responde \texttt{403} en las cinco
          operaciones.
    \item La consulta paginada acepta número de página y tamaño, y devuelve el total de
          elementos junto con la página solicitada.
    \item El conteo por estado usa exclusivamente los códigos de la sección~\ref{sec:estados}.
\end{itemize}
\campo{Método de verificación}{Prueba}
\campo{Estado}{\estadoV}

\subsection{Anteproyecto y tutoría}

\REQ{RF-19}{Carga del anteproyecto en PDF}
\campo{Enunciado}{El sistema debe permitir al estudiante propietario de la solicitud cargar el
anteproyecto en formato PDF, calculando y almacenando su resumen criptográfico SHA-256 en el
momento de la carga.}
\campo{Prioridad}{Must}
\campo{Fuente}{HU-12, CU-12}
\campo{Endpoint}{\cd{POST /api/anteproyectos/enviar/\{solicitudId\}}}
\campo{Rationale}{El anteproyecto es el documento que el tribunal evalúa. Calcular el resumen en
el momento de la carga, y no después, es lo que permite demostrar más tarde que el archivo
entregado es el mismo que se evaluó.}
\campo{Criterio de aceptación}{}
\begin{itemize}
    \item Un archivo cuyo tipo de contenido no sea \texttt{application/pdf} se rechaza con
          \texttt{400} y no se almacena.
    \item Un archivo de más de \textbf{10\,MB} se rechaza con \texttt{400} y no se almacena.
    \item Cargar el anteproyecto de una solicitud ajena responde \texttt{403}.
    \item Si la solicitud está \texttt{SUSPENDIDA}, la carga responde \texttt{400} indicando el
          motivo de suspensión.
    \item Con un anteproyecto ya cargado y no rechazado, un segundo envío responde \texttt{400}:
          el reemplazo exige que el revisor lo rechace primero (RF-21).
    \item El registro almacena el resumen SHA-256 en 64 caracteres hexadecimales y el tamaño en
          bytes.
\end{itemize}
\campo{Método de verificación}{Prueba}
\campo{Estado}{\estadoV}
\campo{Depende de}{RF-13, RF-16}

\REQ{RF-20}{Verificación de integridad del anteproyecto}
\campo{Enunciado}{El sistema debe permitir verificar que el archivo almacenado de un
anteproyecto coincide con el resumen SHA-256 registrado en el momento de su carga, e informar el
resultado de esa comparación.}
\campo{Prioridad}{Should}
\campo{Fuente}{HU-12, CU-12; hallazgo de análisis estático \cd{UNSAFE\_HASH\_EQUALS}
corregido el 2026-08-29}
\campo{Endpoints}{\cd{GET /api/anteproyectos/verificar/\{solicitudId\}},
\cd{GET /api/anteproyectos/ver/\{solicitudId\}},
\cd{GET /api/anteproyectos/solicitud/\{solicitudId\}}}
\campo{Rationale}{Guardar un resumen que nadie comprueba no aporta integridad. La comparación
debe hacerse en tiempo constante para que no filtre información sobre el valor esperado.}
\campo{Criterio de aceptación}{}
\begin{itemize}
    \item Con el archivo íntegro responde \texttt{200} declarando la coincidencia.
    \item Si el archivo fue alterado o falta, responde \texttt{200} declarando explícitamente la
          \textbf{no} coincidencia; no devuelve un error genérico que confunda «alterado» con
          «no encontrado».
    \item La comparación de resúmenes no se realiza con una comparación de cadenas que termine
          en el primer byte distinto.
\end{itemize}
\campo{Método de verificación}{Prueba + Inspección}
\campo{Estado}{\estadoV}

\REQ{RF-21}{Revisión del anteproyecto}
\campo{Enunciado}{El sistema debe permitir a quien posea el permiso
\cd{ANTEPROYECTO\_REVISAR} aprobar o rechazar el anteproyecto cargado, y el rechazo debe
habilitar al estudiante a cargar una versión nueva.}
\campo{Prioridad}{Must}
\campo{Fuente}{HU-12, CU-12}
\campo{Endpoints}{\cd{POST /api/anteproyectos/aprobar/\{id\}},
\cd{POST /api/anteproyectos/rechazar/\{id\}}}
\campo{Rationale}{Es la puerta que RF-19 cierra: el estudiante no puede sustituir el documento
por su cuenta, así que el rechazo del revisor es el único mecanismo que reabre la carga. Sin
este requisito, el bloqueo de RF-19 sería un callejón sin salida.}
\campo{Criterio de aceptación}{}
\begin{itemize}
    \item Aprobar deja el anteproyecto en \texttt{APROBADO}; rechazar, en \texttt{RECHAZADO}.
    \item Tras el rechazo, una nueva carga por RF-19 responde \texttt{200}.
    \item Sin el permiso \cd{ANTEPROYECTO\_REVISAR} responde \texttt{403}.
    \item El estudiante recibe una notificación con el resultado (RF-41).
\end{itemize}
\campo{Método de verificación}{Prueba}
\campo{Estado}{\estadoV}

\REQ{RF-22}{Asignación del tutor}
\campo{Enunciado}{El sistema debe permitir a quien posea el permiso
\cd{TRIBUNAL\_TUTOR\_ASIGNAR} asignar un docente como tutor de una solicitud, y consultar y
retirar esa asignación.}
\campo{Prioridad}{Must}
\campo{Fuente}{HU-03, CU-03}
\campo{Endpoints}{\cd{POST /api/jurados/tutor/asignar}, \cd{POST /api/tutores/asignar},
\cd{GET /api/tutores/solicitud/\{solicitudId\}}, \cd{GET /api/tutores/mis-estudiantes},
\cd{DELETE /api/tutores/\{id\}}}
\campo{Rationale}{El tutor es el primer rol funcional que se fija sobre la solicitud, y de él
depende que puedan abrirse las fases de tutoría (RF-23) y, más adelante, el tribunal (RF-26).}
\campo{Criterio de aceptación}{}
\begin{itemize}
    \item Una solicitud admite un único tutor; asignar un segundo responde \texttt{400}.
    \item La asignación deja la tutoría en estado \texttt{ACTIVO} y la solicitud en
          \texttt{TUTORIA}.
    \item El docente asignado y el estudiante reciben notificación (RF-41).
    \item Sin el permiso \cd{TRIBUNAL\_TUTOR\_ASIGNAR} responde \texttt{403}.
    \item Un docente consulta sus propios estudiantes tutorados sin necesidad de ese permiso, y
          solo los suyos.
\end{itemize}
\campo{Método de verificación}{Prueba}
\campo{Estado}{\estadoV}

\REQ{RF-23}{Gestión de las tres fases de tutoría}
\campo{Enunciado}{El sistema debe permitir al tutor abrir fases de revisión sobre una tutoría,
hasta un máximo de tres, y consultar el resumen y el detalle de esas fases.}
\campo{Prioridad}{Must}
\campo{Fuente}{Auditoría de cobertura del 2026-09-09; reglamento de titulación. Historia de
usuario pendiente.}
\campo{Endpoints}{\cd{POST /api/tutorias/\{tutorId\}/nueva-fase},
\cd{GET /api/tutorias/\{tutorId\}/fases}, \cd{GET /api/tutorias/\{tutorId\}/resumen},
\cd{GET /api/tutorias/estudiante/\{usuarioId\}}, \cd{GET /api/tutorias/docente/\{usuarioId\}}}
\campo{Rationale}{Las tres revisiones son la condición que el reglamento impone antes de
conformar el tribunal. El número está fijado en el sistema y hasta esta versión vivía únicamente
en el código: una regla de negocio que solo existe en el código no es una regla especificada.}
\campo{Criterio de aceptación}{}
\begin{itemize}
    \item El máximo de fases por tutoría es \textbf{3}; crear una cuarta responde \texttt{400}.
    \item Cada fase se numera consecutivamente de 1 a 3 y nace en estado
          \cd{PENDIENTE\_ESTUDIANTE}.
    \item Sin el permiso \cd{TUTORIA\_GESTIONAR} responde \texttt{403}.
    \item Un docente que no es el tutor de esa solicitud responde \texttt{403} aunque posea el
          permiso.
\end{itemize}
\campo{Método de verificación}{Prueba}
\campo{Estado}{\estadoV}

\REQ{RF-24}{Entrega de correcciones del estudiante en una fase}
\campo{Enunciado}{El sistema debe permitir al estudiante cargar un archivo PDF con sus
correcciones en una fase de tutoría abierta, calculando y almacenando su resumen SHA-256.}
\campo{Prioridad}{Must}
\campo{Fuente}{Auditoría de cobertura del 2026-09-09. Historia de usuario pendiente.}
\campo{Endpoints}{\cd{POST /api/tutorias/fases/\{faseId\}/subir-pdf},
\cd{GET /api/tutorias/fases/\{faseId\}/pdf}}
\campo{Rationale}{Es la contraparte del estudiante en el ciclo de tutoría, y la que produce el
documento que en la fase 3 se convierte en el anteproyecto definitivo (RF-26).}
\campo{Criterio de aceptación}{}
\begin{itemize}
    \item Un archivo cuyo tipo de contenido no sea \texttt{application/pdf} se rechaza con
          \texttt{400}.
    \item Cargar en una fase que no esté esperando al estudiante responde \texttt{400}.
    \item Cargar en una fase de otra tutoría responde \texttt{403}.
    \item Sin el permiso \cd{TUTORIA\_AVANCE\_ESTUDIANTE} responde \texttt{403}.
    \item La fase registra el resumen SHA-256 y el tamaño en bytes del archivo entregado.
\end{itemize}
\campo{Método de verificación}{Prueba}
\campo{Estado}{\estadoV}

\REQ{RF-25}{Mensajería de la fase de tutoría}
\campo{Enunciado}{El sistema debe permitir al tutor y al estudiante intercambiar mensajes dentro
de una fase, clasificados como observación, respuesta o aprobación, y marcar los recibidos como
leídos.}
\campo{Prioridad}{Should}
\campo{Fuente}{Auditoría de cobertura del 2026-09-09. Historia de usuario pendiente.}
\campo{Endpoints}{\cd{POST /api/tutorias/fases/\{faseId\}/mensaje},
\cd{PUT /api/tutorias/fases/\{faseId\}/leer}}
\campo{Rationale}{Las observaciones de tutoría deben quedar asociadas a la fase y al documento
que las motivó; fuera del sistema se pierden y no son auditables.}
\campo{Criterio de aceptación}{}
\begin{itemize}
    \item El tipo de mensaje pertenece al conjunto cerrado \texttt{OBSERVACION},
          \texttt{RESPUESTA}, \texttt{APROBACION}; cualquier otro responde \texttt{400}.
    \item Publicar o leer mensajes de una fase ajena responde \texttt{403}.
    \item Un mensaje del tutor de tipo \texttt{OBSERVACION} deja la fase en
          \cd{PENDIENTE\_ESTUDIANTE}.
\end{itemize}
\campo{Método de verificación}{Prueba}
\campo{Estado}{\estadoI}

\REQ{RF-26}{Aprobación de la fase y cierre de la tutoría}
\campo{Enunciado}{El sistema debe permitir al tutor aprobar una fase que tenga entrega del
estudiante y, cuando las tres fases queden aprobadas, cerrar la tutoría promoviendo el documento
de la fase 3 a anteproyecto definitivo con su resumen SHA-256.}
\campo{Prioridad}{Must}
\campo{Fuente}{Auditoría de cobertura del 2026-09-09. Historia de usuario pendiente.}
\campo{Endpoints}{\cd{POST /api/tutorias/fases/\{faseId\}/aprobar},
\cd{POST /api/tutorias/\{tutorId\}/registrar-avance}}
\campo{Rationale}{Este es el punto donde el acompañamiento se convierte en entregable, y la
condición que RF-27 exige para conformar el tribunal. Promover el documento de la fase 3
conservando su resumen evita que el anteproyecto evaluado difiera del último aprobado.}
\campo{Criterio de aceptación}{}
\begin{itemize}
    \item Aprobar una fase sin entrega del estudiante responde \texttt{400}.
    \item Con las tres fases aprobadas, la tutoría queda en \texttt{COMPLETADA}.
    \item Al cerrarse, el anteproyecto de la solicitud queda con el archivo, el resumen SHA-256
          y el tamaño del documento de la fase 3.
    \item Sin el permiso \cd{TUTORIA\_GESTIONAR} responde \texttt{403}; un docente que no es
          el tutor de esa solicitud responde \texttt{403} aunque lo posea.
\end{itemize}
\campo{Método de verificación}{Prueba}
\campo{Estado}{\estadoV}
\campo{Depende de}{RF-23, RF-24}

\subsection{Tribunal, cronograma y salas}

\REQ{RF-27}{Conformación del tribunal}
\campo{Enunciado}{El sistema debe permitir a quien posea el permiso
\cd{TRIBUNAL\_TUTOR\_ASIGNAR} asignar a una solicitud tres docentes con los roles
\texttt{PRESIDENTE}, \cd{VOCAL\_1} y \cd{VOCAL\_2}, solo cuando la tutoría de esa
solicitud esté completada.}
\campo{Prioridad}{Must}
\campo{Fuente}{HU-03, CU-03}
\campo{Endpoints}{\cd{POST /api/jurados/asignar},
\cd{POST /api/jurados/asignar-automatico/\{solicitudId\}},
\cd{GET /api/jurados/solicitud/\{solicitudId\}}, \cd{GET /api/jurados},
\cd{DELETE /api/jurados/\{juradoId\}}}
\campo{Rationale}{Concentra las cuatro reglas que garantizan la validez del tribunal: momento
del proceso, unicidad del docente, unicidad del rol y ausencia de conflicto de interés. Ninguna
de ellas estaba especificada hasta esta versión.}
\campo{Criterio de aceptación}{}
\begin{itemize}
    \item Si la tutoría de la solicitud no está \texttt{COMPLETADA}, responde \texttt{400}
          indicando que faltan las tres revisiones obligatorias.
    \item El rol pertenece al conjunto cerrado \texttt{PRESIDENTE}, \cd{VOCAL\_1},
          \cd{VOCAL\_2}; cualquier otro responde \texttt{400}.
    \item Un docente ya asignado a esa solicitud no puede asignarse de nuevo: responde
          \texttt{400}.
    \item Un rol ya ocupado en esa solicitud no puede volver a asignarse: responde \texttt{400}.
    \item \textbf{El tutor de la solicitud no puede ser además jurado de la misma solicitud};
          el intento responde \texttt{400} por conflicto de interés.
    \item Un docente marcado como no disponible no puede asignarse: responde \texttt{400}.
    \item Conformado el tribunal, la solicitud pasa a \texttt{EVALUACION}.
\end{itemize}
\campo{Método de verificación}{Prueba}
\campo{Estado}{\estadoV}
\campo{Depende de}{RF-26}

\REQ{RF-28}{Asignación masiva de tribunales y sugerencia de jurados}
\campo{Enunciado}{El sistema debe permitir conformar tribunales para varias solicitudes en una
sola operación, y sugerir docentes candidatos para una solicitud dada.}
\campo{Prioridad}{Should}
\campo{Fuente}{HU-03, CU-03; ADR-006 (la asignación masiva se resuelve en procedimiento
almacenado)}
\campo{Endpoints}{\cd{POST /api/jurados/asignar-masivo},
\cd{GET /api/jurados/sugerencias/\{solicitudId\}},
\cd{GET /api/jurados/docente/\{docenteId\}},
\cd{GET /api/jurados/tutor/docente/\{docenteId\}},
\cd{GET /api/jurados/info/\{solicitudId\}/\{usuarioId\}}}
\campo{Rationale}{En una convocatoria la coordinación conforma decenas de tribunales el mismo
día; hacerlo de uno en uno es donde aparecen los errores de rol repetido y de conflicto de
interés que RF-27 rechaza.}
\campo{Criterio de aceptación}{}
\begin{itemize}
    \item La operación masiva aplica a cada solicitud las mismas cinco reglas de RF-27, y
          devuelve por cada una si se asignó o el motivo del rechazo, sin abortar el lote
          completo por un fallo individual.
    \item La sugerencia nunca propone al tutor de esa solicitud ni a un docente ya asignado.
    \item Un docente puede consultar sus propias asignaciones como jurado o tutor sin el permiso
          \cd{TRIBUNAL\_TUTOR\_ASIGNAR}, y solo las suyas.
\end{itemize}
\campo{Método de verificación}{Prueba}
\campo{Estado}{\estadoV}

\REQ{RF-29}{Programación del cronograma de defensa}
\campo{Enunciado}{El sistema debe programar la defensa de una solicitud en una sala, fecha y
hora determinadas, con una duración de 45 minutos, rechazando la programación si la sala o
alguno de los miembros del tribunal ya está ocupado en ese intervalo.}
\campo{Prioridad}{Must}
\campo{Fuente}{HU-04, CU-04; ADR-006 (la validación de conflicto de jurado es un procedimiento
almacenado)}
\campo{Endpoints}{\cd{POST /api/cronogramas/crear},
\cd{POST /api/cronogramas/auto/\{solicitudId\}}, \cd{GET /api/cronogramas},
\cd{GET /api/cronogramas/solicitud/\{id\}}, \cd{GET /api/cronogramas/estudiante/\{id\}},
\cd{GET /api/cronogramas/usuario/\{id\}}, \cd{DELETE /api/cronogramas/\{id\}}}
\campo{Rationale}{Un choque de agenda descubierto el día de la defensa cuesta una reprogramación
completa. La duración fija de 45 minutos y la ventana de búsqueda automática de 30 días son
reglas de negocio que hasta esta versión vivían solo en el código.}
\campo{Criterio de aceptación}{}
\begin{itemize}
    \item La duración de toda defensa programada es de \textbf{45 minutos}.
    \item Si la sala tiene otra defensa que se solapa con ese intervalo, responde \texttt{400}
          nombrando la sala.
    \item Si algún miembro del tribunal tiene otra defensa que se solapa, responde \texttt{400}
          nombrando al docente.
    \item La programación automática busca la primera franja libre en los \textbf{próximos 30
          días} sobre las salas disponibles; si no encuentra ninguna responde \texttt{400}
          declarándolo, en vez de programar en una franja ocupada.
    \item Sin el permiso \cd{CRONOGRAMA\_GESTIONAR} responde \texttt{403} en la creación y
          en la eliminación.
    \item El cronograma creado queda en estado \texttt{PROGRAMADO}.
\end{itemize}
\campo{Método de verificación}{Prueba}
\campo{Estado}{\estadoV}
\campo{Depende de}{RF-27, RF-31}

\REQ{RF-30}{Consulta de disponibilidad de salas y franjas}
\campo{Enunciado}{El sistema debe informar las franjas horarias disponibles y verificar la
disponibilidad de una sala en un instante dado, antes de intentar la programación.}
\campo{Prioridad}{Should}
\campo{Fuente}{HU-04, HU-10}
\campo{Endpoints}{\cd{GET /api/cronogramas/disponibilidad},
\cd{GET /api/cronogramas/verificar-disponibilidad}}
\campo{Rationale}{Sin consulta previa, la única forma de saber si una franja está libre es
intentar programar y recibir el rechazo de RF-29, lo que convierte una consulta en un intento
fallido.}
\campo{Criterio de aceptación}{}
\begin{itemize}
    \item La disponibilidad devuelta es coherente con las reglas de conflicto de RF-29: una
          franja informada como libre no puede ser rechazada por solapamiento de sala.
    \item La verificación puntual responde \texttt{200} declarando disponible o no disponible,
          sin crear ningún cronograma.
\end{itemize}
\campo{Método de verificación}{Prueba}
\campo{Estado}{\estadoI}

\REQ{RF-31}{Gestión de salas de sustentación}
\campo{Enunciado}{El sistema debe permitir consultar las salas y, a quien posea el permiso
\cd{SALA\_GESTIONAR}, crearlas y eliminarlas, con código único, nombre, capacidad y
disponibilidad.}
\campo{Prioridad}{Should}
\campo{Fuente}{HU-10, CU-10}
\campo{Endpoints}{\cd{GET /api/salas}, \cd{GET /api/salas/paginado}, \cd{POST /api/salas},
\cd{DELETE /api/salas/\{id\}}}
\campo{Rationale}{Las salas son el recurso escaso que RF-29 reparte; sin catálogo, la
programación automática no tiene sobre qué buscar.}
\campo{Criterio de aceptación}{}
\begin{itemize}
    \item Un código de sala duplicado responde \texttt{400}.
    \item Eliminar una sala con defensas programadas responde \texttt{400} en vez de dejar
          cronogramas huérfanos.
    \item Sin el permiso \cd{SALA\_GESTIONAR} la creación y la eliminación responden
          \texttt{403}; la consulta está abierta a cualquier usuario autenticado.
    \item Una sala marcada como no disponible no es candidata de la programación automática de
          RF-29.
\end{itemize}
\campo{Método de verificación}{Prueba}
\campo{Estado}{\estadoI}
\campo{Nota}{Era el requisito RF-10 de la especificación anterior, declarado allí «no verificado» porque
la clase de prueba que citaba no existía. Sigue sin prueba dedicada; el estado se declara
\estadoI{} conforme al vocabulario de la sección~\ref{sec:estados-req}, que distingue construido
de verificado.}

\subsection{Evaluación por rúbrica}

\REQ{RF-32}{Gestión de rúbricas y de sus criterios}
\campo{Enunciado}{El sistema debe permitir a quien posea el permiso
\cd{RUBRICA\_GESTIONAR} crear y eliminar rúbricas, definir sus criterios y sembrar el
conjunto de criterios por omisión de una rúbrica, y debe permitir a cualquier usuario
autenticado consultarlas.}
\campo{Prioridad}{Must}
\campo{Fuente}{HU-05, CU-05}
\campo{Endpoints}{\cd{GET /api/rubricas}, \cd{GET /api/rubricas/\{id\}},
\cd{POST /api/rubricas}, \cd{DELETE /api/rubricas/\{id\}},
\cd{POST /api/rubricas/\{rubricaId\}/criterios},
\cd{GET /api/rubricas/\{rubricaId\}/criterios},
\cd{POST /api/rubricas/\{rubricaId\}/inicializar-criterios}}
\campo{Rationale}{La rúbrica es el instrumento que hace comparable la evaluación entre
tribunales distintos. Sin criterios definidos, RF-33 no tiene sobre qué registrar.}
\campo{Criterio de aceptación}{}
\begin{itemize}
    \item Sin el permiso \cd{RUBRICA\_GESTIONAR}, crear, eliminar o definir criterios
          responde \texttt{403}; la consulta está abierta a cualquier usuario autenticado.
    \item Eliminar una rúbrica ya utilizada en una evaluación responde \texttt{400}.
    \item La siembra de criterios por omisión es idempotente: repetirla no duplica criterios.
\end{itemize}
\campo{Método de verificación}{Prueba}
\campo{Estado}{\estadoV}

\REQ{RF-33}{Registro de la evaluación por criterio de rúbrica}
\campo{Enunciado}{El sistema debe permitir a un miembro del tribunal registrar su calificación
para \textbf{todos} los criterios de la rúbrica de una solicitud, en una escala de 1 a 100 por
criterio, y solo a nombre propio.}
\campo{Prioridad}{Must}
\campo{Fuente}{HU-05, CU-05; hallazgo de auditoría OWASP A01 (escritura a nombre de otro
jurado)}
\campo{Endpoints}{\cd{POST /api/rubrica-evaluacion/registrar},
\cd{POST /api/evaluaciones-jurado/guardar},
\cd{GET /api/rubrica-evaluacion/criterios/\{rubricaId\}},
\cd{GET /api/rubrica-evaluacion/solicitud/\{solicitudId\}/jurado/\{juradoId\}}}
\campo{Rationale}{Una evaluación parcial no es comparable ni promediable, y una evaluación
registrada a nombre de otro jurado invalida el acta. Ambas condiciones tienen que rechazarse en
el servidor, no confiarse al cliente.}
\campo{Criterio de aceptación}{}
\begin{itemize}
    \item Una escala fuera del intervalo \textbf{1 a 100} responde \texttt{400} declarando el
          valor recibido.
    \item Una entrega que no cubra \textbf{todos} los criterios de la rúbrica responde
          \texttt{400} indicando cuántos criterios exige la rúbrica.
    \item Registrar la evaluación a nombre de otro jurado responde \texttt{403}, aunque quien lo
          intente posea \cd{EVALUACION\_RUBRICA\_REGISTRAR}.
    \item Registrar una evaluación en una solicitud en la que quien la envía no es jurado
          responde \texttt{400}.
    \item Si la rúbrica no tiene criterios definidos, responde \texttt{400}.
\end{itemize}
\campo{Método de verificación}{Prueba}
\campo{Estado}{\estadoV}
\campo{Depende de}{RF-27, RF-32, RNF-13}

\REQ{RF-34}{Cálculo de la nota final ponderada}
\campo{Enunciado}{El sistema debe calcular la nota final de una solicitud como la ponderación de
la nota del instructor y la nota promedio del tribunal, con pesos que sumen 100, y debe derivar
de esa nota el resultado \texttt{APROBADO} o \texttt{REPROBADO}.}
\campo{Prioridad}{Must}
\campo{Fuente}{HU-05, CU-05; ADR-006 (el promedio se resuelve en procedimiento almacenado)}
\campo{Endpoints}{\cd{POST /api/evaluaciones/evaluar-ponderado},
\cd{POST /api/evaluaciones/evaluar},
\cd{POST /api/evaluaciones/calcular-promedio/\{solicitudId\}}}
\campo{Rationale}{Es la regla que decide si un estudiante aprueba. Sus tres constantes ---los
pesos por omisión, el rango de nota y el umbral de aprobación--- vivían únicamente en el código
hasta esta versión, de modo que ni el estudiante ni el tribunal podían conocerlas leyendo la
especificación.}
\campo{Criterio de aceptación}{}
\begin{itemize}
    \item Los pesos por omisión son \textbf{60\,\% instructor} y \textbf{40\,\% tribunal}.
    \item Si los pesos enviados no suman \textbf{100} (con tolerancia $0{,}01$), responde
          \texttt{400} indicando la suma recibida y no calcula nada.
    \item Toda nota debe estar en el intervalo \textbf{0 a 10}; fuera de él responde
          \texttt{400}.
    \item El resultado es \texttt{APROBADO} si la nota final es \textbf{mayor o igual a 7,0}, y
          \texttt{REPROBADO} en otro caso. El resultado es calculado y no admite fijarse
          directamente.
    \item Registrada la nota final, la solicitud pasa a \texttt{CALIFICADA}.
    \item Sin el permiso \cd{EVALUACION\_CALIFICAR} responde \texttt{403}.
    \item El estudiante recibe notificación de su nota final (RF-41).
\end{itemize}
\campo{Método de verificación}{Prueba}
\campo{Estado}{\estadoV}
\campo{Depende de}{RF-33}

\REQ{RF-35}{Consulta de resultados de evaluación}
\campo{Enunciado}{El sistema debe permitir consultar la evaluación de una solicitud, la nota del
tribunal, las evaluaciones de un estudiante y las observaciones registradas, restringiendo cada
consulta a quien tenga derecho sobre ese expediente.}
\campo{Prioridad}{Should}
\campo{Fuente}{HU-05, CU-05}
\campo{Endpoints}{\cd{GET /api/evaluaciones}, \cd{GET /api/evaluaciones/solicitud/\{solicitudId\}},
\cd{GET /api/evaluaciones/estudiante/\{estudianteId\}},
\cd{GET /api/evaluaciones/usuario/\{usuarioId\}},
\cd{GET /api/rubrica-evaluacion/solicitud/\{solicitudId\}},
\cd{GET /api/rubrica-evaluacion/nota-tribunal/\{solicitudId\}},
\cd{GET /api/rubrica-evaluacion/observaciones/\{solicitudId\}},
\cd{GET /api/evaluaciones-jurado/tribunal/\{solicitudId\}},
\cd{GET /api/evaluaciones-jurado/\{solicitudId\}/\{juradoId\}}}
\campo{Rationale}{La calificación es un dato personal sensible del estudiante; el acceso no
puede depender de conocer el identificador del expediente.}
\campo{Criterio de aceptación}{}
\begin{itemize}
    \item Consultar las evaluaciones de otro estudiante o de otro usuario responde \texttt{403}.
    \item El listado general de evaluaciones exige el permiso \cd{EVALUACION\_CALIFICAR};
          sin él responde \texttt{403}.
    \item La consulta de una solicitud sin evaluación registrada responde \texttt{404}, no una
          nota vacía que pueda confundirse con un cero.
\end{itemize}
\campo{Método de verificación}{Prueba}
\campo{Estado}{\estadoV}
\campo{Depende de}{RNF-13}

\subsection{Acta de pre-sustentación}

\REQ{RF-36}{Generación del acta}
\campo{Enunciado}{El sistema debe generar el acta de pre-sustentación de una solicitud en
formato PDF, con los datos del estudiante, la composición del tribunal, la evaluación y el
resultado, y dejarla en estado \texttt{GENERADA}.}
\campo{Prioridad}{Must}
\campo{Fuente}{HU-06, CU-06}
\campo{Endpoint}{\cd{POST /api/v1/actas/generar/\{solicitudId\}}}
\campo{Rationale}{El acta es el documento con valor académico que sale del sistema. Exigir que
exista evaluación final antes de emitirla impide producir un acta sin resultado, que sería un
documento oficial vacío.}
\campo{Criterio de aceptación}{}
\begin{itemize}
    \item Generar el acta de una solicitud sin evaluación final registrada responde
          \texttt{400}.
    \item El acta generada queda en estado \texttt{GENERADA} y con su archivo PDF disponible.
    \item El PDF contiene los datos del estudiante, la tabla del tribunal con el rol de cada
          miembro y su confirmación, y la evaluación con el resultado.
    \item Sin el permiso \cd{ACTA\_GENERAR} responde \texttt{403}.
    \item La generación queda registrada en el historial de estados del acta (RF-40).
\end{itemize}
\campo{Método de verificación}{Prueba}
\campo{Estado}{\estadoV}
\campo{Depende de}{RF-34}

\REQ{RF-37}{Firma digital multi-actor del acta}
\campo{Enunciado}{El sistema debe permitir que el acta sea firmada por los cuatro roles
\texttt{PRESIDENTE}, \cd{VOCAL\_1}, \cd{VOCAL\_2} y \texttt{TUTOR}, aceptando cada firma
únicamente del docente que ejerce ese rol en esa solicitud, y debe cerrar el proceso cuando las
cuatro firmas estén aplicadas.}
\campo{Prioridad}{Must}
\campo{Fuente}{HU-07, CU-07; ADR-006 (la firma se persiste en procedimiento almacenado)}
\campo{Endpoint}{\cd{POST /api/v1/actas/firmar/\{actaId\}}}
\campo{Rationale}{Es el requisito que cierra todo el proceso, y el que mejor ilustra por qué un
permiso no basta como control de acceso: poseer \cd{ACTA\_FIRMAR} habilita la funcionalidad,
pero firmar como \texttt{PRESIDENTE} exige ser el presidente de \emph{ese} tribunal.}
\campo{Criterio de aceptación}{}
\begin{itemize}
    \item El rol de firma pertenece al conjunto cerrado \texttt{PRESIDENTE}, \cd{VOCAL\_1},
          \cd{VOCAL\_2}, \texttt{TUTOR}; cualquier otro responde \texttt{400}.
    \item Firmar con un rol que el firmante no ejerce en esa solicitud responde \texttt{403},
          aunque posea \cd{ACTA\_FIRMAR}.
    \item Cada firma registra el rol, la fecha y la observación asociada.
    \item Aplicada la cuarta firma, el acta pasa a \texttt{FINALIZADA}, la solicitud pasa a
          \texttt{COMPLETADA} y el PDF se regenera reflejando el estado final de las firmas.
    \item El estudiante recibe notificación del cierre del proceso (RF-41).
\end{itemize}
\campo{Método de verificación}{Prueba}
\campo{Estado}{\estadoV}
\campo{Depende de}{RF-36, RF-27, RF-22}
\campo{Nota}{Esta versión eleva el requisito de \emph{Should} a \textbf{Must}. La razón es que
RF-37 es la única transición que lleva la solicitud a \texttt{COMPLETADA}: sin él, ningún
expediente puede cerrarse, y un requisito del que depende el cierre del proceso no puede ser
opcional. La especificación anterior lo declaraba \emph{Should} sin justificar esa
prioridad.}

\REQ{RF-38}{Consulta y descarga del acta}
\campo{Enunciado}{El sistema debe permitir consultar y descargar el acta a los actores con
derecho sobre ese expediente, y permitir a un docente listar exclusivamente las actas de las
pre-sustentaciones en las que participa.}
\campo{Prioridad}{Should}
\campo{Fuente}{Auditoría de cobertura del 2026-09-09. Historia de usuario pendiente.}
\campo{Endpoints}{\cd{GET /api/v1/actas}, \cd{GET /api/v1/actas/mis-actas},
\cd{GET /api/v1/actas/buscar}, \cd{GET /api/v1/actas/\{id\}},
\cd{GET /api/v1/actas/solicitud/\{solicitudId\}},
\cd{GET /api/v1/actas/ver/\{actaId\}}, \cd{GET /api/v1/actas/descargar/\{actaId\}},
\cd{DELETE /api/v1/actas/\{id\}}}
\campo{Rationale}{Separar el listado completo (\cd{ACTAS\_VER}) del listado propio
(\cd{ACTAS\_VER\_PROPIAS}) es lo que permite dar acceso al docente sin exponerle el
expediente de toda la facultad.}
\campo{Criterio de aceptación}{}
\begin{itemize}
    \item El listado propio devuelve únicamente las actas en las que quien consulta es tutor o
          miembro del tribunal.
    \item Consultar o descargar un acta sobre la que no se tiene derecho responde \texttt{403},
          conociendo su identificador.
    \item Descargar un acta cuyo PDF no se ha generado responde \texttt{400}.
    \item La eliminación exige el permiso \cd{ACTAS\_GESTIONAR}; sin él responde
          \texttt{403}.
\end{itemize}
\campo{Método de verificación}{Prueba + Demostración}
\campo{Estado}{\estadoV}
\campo{Depende de}{RNF-13}

\REQ{RF-39}{Cambio de estado del acta}
\campo{Enunciado}{El sistema debe permitir cambiar el estado de un acta únicamente por
transiciones válidas conforme a la sección~\ref{sec:estados}, exigiendo motivo para
\texttt{OBSERVADA} y \texttt{ANULADA}.}
\campo{Prioridad}{Should}
\campo{Fuente}{Auditoría de cobertura del 2026-09-09. Historia de usuario pendiente.}
\campo{Endpoint}{\cd{PATCH /api/v1/actas/\{id\}/estado}}
\campo{Rationale}{Sin una máquina de estados aplicada en el servidor, un acta podría pasar de
\texttt{FINALIZADA} a \texttt{REVISADA} y reabrir un expediente cerrado sin dejar rastro de por
qué.}
\campo{Criterio de aceptación}{}
\begin{itemize}
    \item Una transición no permitida responde \texttt{400} nombrando el estado de origen, el de
          destino y los estados válidos.
    \item Pasar a \texttt{OBSERVADA} o a \texttt{ANULADA} sin motivo responde \texttt{400}.
    \item Fijar el estado que el acta ya tiene responde \texttt{400}.
    \item La transición a \texttt{ANULADA} desde cualquier estado está reservada al
          Administrador; para el resto de actores rigen las transiciones de la tabla.
    \item Sin el permiso \cd{ACTA\_ESTADO\_CAMBIAR} responde \texttt{403}.
\end{itemize}
\campo{Método de verificación}{Prueba + Demostración}
\campo{Estado}{\estadoV}

\REQ{RF-40}{Historial de estados del acta}
\campo{Enunciado}{El sistema debe registrar cada cambio de estado de un acta con el usuario, su
rol, el estado anterior, el estado nuevo, el motivo y la fecha, y permitir consultar ese
historial a quien tenga derecho sobre el acta.}
\campo{Prioridad}{Should}
\campo{Fuente}{Auditoría de cobertura del 2026-09-09. Historia de usuario pendiente.}
\campo{Endpoint}{\cd{GET /api/v1/actas/\{id\}/historial}}
\campo{Rationale}{El acta es el documento con efectos académicos; saber quién la observó o la
anuló, y por qué, es parte del documento tanto como su contenido.}
\campo{Criterio de aceptación}{}
\begin{itemize}
    \item Cada cambio de estado de RF-39 y la generación de RF-36 producen exactamente una
          entrada de historial.
    \item Consultar el historial de un acta ajena responde \texttt{403}.
    \item Sin el permiso \cd{ACTA\_HISTORIAL\_VER} responde \texttt{403}.
    \item Las entradas del historial no son editables ni eliminables por la API.
\end{itemize}
\campo{Método de verificación}{Prueba + Inspección}
\campo{Estado}{\estadoV}

\subsection{Comunicación y apoyo al usuario}

\REQ{RF-41}{Notificaciones al usuario}
\campo{Enunciado}{El sistema debe registrar una notificación dirigida al usuario afectado en
cada hito del proceso, y permitirle consultarlas, filtrar las no leídas, marcarlas como leídas
individualmente o en bloque, y eliminarlas.}
\campo{Prioridad}{Should}
\campo{Fuente}{HU-08, CU-08}
\campo{Endpoints}{\cd{GET /api/v1/notificaciones/usuario/\{usuarioId\}},
\cd{GET /api/v1/notificaciones/usuario/\{usuarioId\}/no-leidas},
\cd{PATCH /api/v1/notificaciones/\{id\}/marcar-leida},
\cd{PATCH /api/v1/notificaciones/usuario/\{usuarioId\}/marcar-todas-leidas},
\cd{DELETE /api/v1/notificaciones/\{id\}}}
\campo{Rationale}{Los hitos del proceso ocurren sin que el afectado esté conectado. La
notificación interna es lo que evita que el estudiante tenga que consultar su estado a diario
para enterarse de una decisión.}
\campo{Criterio de aceptación}{}
\begin{itemize}
    \item Se registra notificación al menos en: aprobación o rechazo de la solicitud (RF-15),
          resultado de la revisión del anteproyecto (RF-21), asignación de tutor (RF-22), carga
          del anteproyecto, nota final (RF-34) y cierre del acta (RF-37).
    \item Consultar, marcar o eliminar notificaciones de otro usuario responde \texttt{403}.
    \item El listado global de notificaciones del sistema exige el permiso
          \cd{NOTIFICACIONES\_GLOBAL\_VER}; sin él responde \texttt{403}.
    \item Si el envío de correo saliente está desactivado o el proveedor falla, la notificación
          interna \textbf{se registra igualmente}: la entrega por correo es un canal adicional,
          no la notificación.
\end{itemize}
\campo{Método de verificación}{Prueba}
\campo{Estado}{\estadoV}

\REQ{RF-42}{Envío manual de notificaciones}
\campo{Enunciado}{El sistema debe permitir a quien posea el permiso
\cd{NOTIFICACIONES\_ENVIAR} crear una notificación dirigida a un usuario determinado.}
\campo{Prioridad}{Could}
\campo{Fuente}{HU-08}
\campo{Endpoints}{\cd{POST /api/v1/notificaciones/crear}, \cd{GET /api/v1/notificaciones}}
\campo{Rationale}{Cubre las comunicaciones que no corresponden a un hito automatizado, sin
obligar a la coordinación a salir del sistema.}
\campo{Criterio de aceptación}{}
\begin{itemize}
    \item Sin el permiso \cd{NOTIFICACIONES\_ENVIAR} responde \texttt{403}.
    \item El destinatario debe existir; en otro caso responde \texttt{400}.
    \item La notificación creada aparece en el listado del destinatario (RF-41).
\end{itemize}
\campo{Método de verificación}{Prueba}
\campo{Estado}{\estadoI}

\REQ{RF-43}{Consulta del estado del proceso en tiempo real}
\campo{Enunciado}{El sistema debe exponer, en una sola petición, el estado consolidado del
proceso de una solicitud, y admitir la consulta consolidada de varias solicitudes en lote.}
\campo{Prioridad}{Should}
\campo{Fuente}{Auditoría de cobertura del 2026-09-09. Historia de usuario pendiente.}
\campo{Endpoints}{\cd{GET /api/estado/solicitud/\{id\}},
\cd{POST /api/estado/solicitudes/batch}}
\campo{Rationale}{El cliente refresca el estado cada 15 segundos. Sin un endpoint consolidado
tendría que encadenar cinco llamadas por solicitud, multiplicando por cinco la carga de un
sondeo que ya es periódico. La actualización se resuelve por sondeo y \textbf{no} por
\emph{WebSocket}; declararlo evita que la política de seguridad de contenido reserve permisos
para un canal que no existe.}
\campo{Criterio de aceptación}{}
\begin{itemize}
    \item Devuelve en una respuesta el estado de la solicitud, la tutoría, el tribunal, el
          cronograma y el acta.
    \item La consulta en lote devuelve un elemento por identificador solicitado, sin fallar el
          lote completo por una solicitud inexistente.
    \item Consultar el estado de una solicitud ajena responde \texttt{403}.
\end{itemize}
\campo{Método de verificación}{Prueba}
\campo{Estado}{\estadoI}

\REQ{RF-44}{Asistente virtual de ayuda}
\campo{Enunciado}{El sistema debe responder consultas de ayuda del usuario autenticado sobre los
módulos del sistema, y ofrecer el conjunto de temas disponibles en cada respuesta.}
\campo{Prioridad}{Could}
\campo{Fuente}{Auditoría de cobertura del 2026-09-09. Historia de usuario pendiente.}
\campo{Endpoint}{\cd{POST /api/v1/chatbot/ask}}
\campo{Rationale}{Reduce las consultas de trámite a la coordinación. La resolución es interna y
determinista, sin consumir ningún proveedor externo de modelo de lenguaje: eso mantiene el
requisito dentro del alcance de RNF-20 y evita que el mensaje del usuario salga del sistema.}
\campo{Criterio de aceptación}{}
\begin{itemize}
    \item Sin autenticación responde indicando que se requiere iniciar sesión, sin procesar el
          mensaje.
    \item Ante un mensaje que no reconoce, responde declarando que no puede responderlo y
          ofrece los temas disponibles, en vez de inventar una respuesta.
    \item \textbf{El mensaje del usuario no se envía a ningún servicio externo}, y ninguna
          respuesta incluye datos de otro usuario.
    \item Toda respuesta incluye la lista de temas soportados.
\end{itemize}
\campo{Método de verificación}{Prueba + Inspección}
\campo{Estado}{\estadoV}
\campo{Nota}{El asistente ofrece hoy orientación sobre cambio y recuperación de contraseña,
capacidades que no existen (RF-05, RF-06). Es una promesa que el sistema no puede cumplir, y se
resuelve implementando esos requisitos o retirando ese tema del asistente; no dejando la
contradicción en pie.}

\REQ{RF-45}{Consulta del directorio de universidades}
\campo{Enunciado}{El sistema debe ofrecer al usuario autenticado el listado de universidades del
Ecuador obtenido del directorio externo, con resultado cacheado y con una respuesta de reserva
cuando el proveedor no está disponible.}
\campo{Prioridad}{Could}
\campo{Fuente}{Criterio de consumo de API externa de la asignatura}
\campo{Endpoint}{\cd{GET /api/universidades}}
\campo{Rationale}{Es la única dependencia externa de datos del sistema. La reserva local existe
para que la indisponibilidad de un tercero no se convierta en un error de esta aplicación.}
\campo{Criterio de aceptación}{}
\begin{itemize}
    \item Con el proveedor disponible responde \texttt{200} con el listado.
    \item El resultado se sirve desde caché durante \textbf{10 minutos} antes de volver a
          consultar al proveedor.
    \item Ante error o tiempo de espera del proveedor, tras \textbf{3 reintentos} con retroceso
          exponencial desde 1\,s, responde \texttt{200} con la entrada de reserva local; nunca
          propaga un \texttt{5xx} al cliente.
    \item Los tiempos de espera de conexión, lectura y escritura son de \textbf{5 segundos}.
    \item \textbf{Ninguna petición al proveedor incluye datos de un usuario del sistema.}
\end{itemize}
\campo{Método de verificación}{Prueba + Análisis}
\campo{Estado}{\estadoV}
\campo{Nota}{El análisis estadístico de caché fría contra caché caliente sobre este endpoint
está archivado en \cd{k6/README.md} (30 muestras por escenario, prueba de Wilcoxon, IC 95\,\%).}

\subsection{Centro de orientación de titulación}

\REQ{RF-46}{Exploración del catálogo de temas propuestos}
\campo{Enunciado}{El sistema debe permitir a quien posea el permiso
\cd{ORIENTACION\_TEMAS\_VER} explorar el catálogo de temas de titulación propuestos y
consultar el detalle de un tema.}
\campo{Prioridad}{Should}
\campo{Fuente}{\cd{docs/funcionalidades/CENTRO-ORIENTACION.md}. Historia de usuario pendiente.}
\campo{Endpoints}{\cd{GET /api/v1/orientacion/temas},
\cd{GET /api/v1/orientacion/temas/\{temaId\}}}
\campo{Rationale}{La elección del tema es el punto donde más estudiantes se detienen. El
catálogo convierte esa decisión en una elección informada sobre líneas de investigación reales
de la carrera.}
\campo{Criterio de aceptación}{}
\begin{itemize}
    \item Los cuatro roles poseen \cd{ORIENTACION\_TEMAS\_VER} en la asignación inicial
          (sección~\ref{sec:permisos}); sin el permiso responde \texttt{403}.
    \item El catálogo admite filtrado por línea de investigación y por área temática.
    \item Un tema inexistente responde \texttt{404}.
\end{itemize}
\campo{Método de verificación}{Prueba}
\campo{Estado}{\estadoV}

\REQ{RF-47}{Generación de ideas de tema}
\campo{Enunciado}{El sistema debe generar propuestas de tema a partir de la línea de
investigación y el área temática indicadas por el usuario, sin persistirlas automáticamente.}
\campo{Prioridad}{Could}
\campo{Fuente}{\cd{docs/funcionalidades/CENTRO-ORIENTACION.md}. Historia de usuario pendiente.}
\campo{Endpoint}{\cd{POST /api/v1/orientacion/temas/generar}}
\campo{Rationale}{Una sugerencia es una recomendación, no un hecho: persistirla sola convertiría
una propuesta que nadie eligió en parte del expediente. Solo el acto explícito de RF-48 la
guarda. La generación es \textbf{interna}, sin proveedor externo, en coherencia con RNF-20.}
\campo{Criterio de aceptación}{}
\begin{itemize}
    \item Si el área temática indicada no pertenece a la línea de investigación declarada,
          responde \texttt{400}.
    \item Las propuestas devueltas \textbf{no se persisten} como temas del estudiante.
    \item \textbf{Ningún dato del usuario se envía a un servicio externo} en esta operación.
    \item Sin el permiso \cd{ORIENTACION\_TEMAS\_VER} responde \texttt{403}.
\end{itemize}
\campo{Método de verificación}{Prueba + Inspección}
\campo{Estado}{\estadoI}

\REQ{RF-48}{Temas guardados por el estudiante}
\campo{Enunciado}{El sistema debe permitir a un usuario con rol \texttt{ESTUDIANTE} guardar y
retirar temas de su lista personal, y consultarla.}
\campo{Prioridad}{Could}
\campo{Fuente}{\cd{docs/funcionalidades/CENTRO-ORIENTACION.md}. Historia de usuario pendiente.}
\campo{Endpoints}{\cd{GET /api/v1/orientacion/temas/guardados},
\cd{POST /api/v1/orientacion/temas/\{temaId\}/guardar},
\cd{DELETE /api/v1/orientacion/temas/\{temaId\}/guardar}}
\campo{Rationale}{Es lo que hace del catálogo una herramienta de decisión y no solo un listado:
permite comparar candidatos antes de comprometerse en RF-13.}
\campo{Criterio de aceptación}{}
\begin{itemize}
    \item Guardar un tema ya guardado responde \texttt{400} sin duplicar la entrada.
    \item Retirar un tema que no estaba guardado responde \texttt{400}.
    \item La lista devuelta contiene exclusivamente los temas del estudiante en sesión.
    \item Un usuario cuyo rol no es \texttt{ESTUDIANTE} responde \texttt{403}.
\end{itemize}
\campo{Método de verificación}{Prueba}
\campo{Estado}{\estadoV}

\REQ{RF-49}{Gestión del catálogo de orientación}
\campo{Enunciado}{El sistema debe permitir a quien posea el permiso
\cd{ORIENTACION\_CATALOGO\_GESTIONAR} crear, editar y eliminar temas propuestos y recursos
de titulación.}
\campo{Prioridad}{Should}
\campo{Fuente}{\cd{docs/funcionalidades/CENTRO-ORIENTACION.md}. Historia de usuario pendiente.}
\campo{Endpoints}{\cd{POST /api/v1/orientacion/temas},
\cd{PUT /api/v1/orientacion/temas/\{temaId\}},
\cd{DELETE /api/v1/orientacion/temas/\{temaId\}},
\cd{POST /api/v1/orientacion/recursos},
\cd{PUT /api/v1/orientacion/recursos/\{id\}},
\cd{DELETE /api/v1/orientacion/recursos/\{id\}}}
\campo{Rationale}{RF-46 da por hecho un catálogo curado; sin este requisito no existiría el
mecanismo que lo mantiene, y el catálogo quedaría congelado en su siembra inicial.}
\campo{Criterio de aceptación}{}
\begin{itemize}
    \item Sin el permiso \cd{ORIENTACION\_CATALOGO\_GESTIONAR} las seis operaciones
          responden \texttt{403}.
    \item Un tema con el área temática incoherente con su línea de investigación responde
          \texttt{400}.
    \item Eliminar un tema guardado por algún estudiante lo retira también de esas listas
          personales, sin dejar referencias colgantes.
\end{itemize}
\campo{Método de verificación}{Prueba}
\campo{Estado}{\estadoV}

\REQ{RF-50}{Consulta de recursos de titulación}
\campo{Enunciado}{El sistema debe ofrecer a cualquier usuario autenticado el catálogo de
recursos de apoyo a la titulación.}
\campo{Prioridad}{Could}
\campo{Fuente}{\cd{docs/funcionalidades/CENTRO-ORIENTACION.md}. Historia de usuario pendiente.}
\campo{Endpoint}{\cd{GET /api/v1/orientacion/recursos}}
\campo{Rationale}{Los recursos son material de referencia; restringirlos por permiso no aportaría
protección y sí fricción.}
\campo{Criterio de aceptación}{}
\begin{itemize}
    \item Responde \texttt{200} con los recursos vigentes a cualquier usuario autenticado.
    \item Sin autenticación responde \texttt{401}.
\end{itemize}
\campo{Método de verificación}{Prueba}
\campo{Estado}{\estadoV}

\REQ{RF-51}{Ruta de progreso de titulación}
\campo{Enunciado}{El sistema debe presentar al estudiante los ocho pasos de la ruta de
titulación, permitirle marcar los completados, y calcular el porcentaje de avance.}
\campo{Prioridad}{Should}
\campo{Fuente}{\cd{docs/funcionalidades/CENTRO-ORIENTACION.md}. Historia de usuario pendiente.}
\campo{Endpoints}{\cd{GET /api/v1/orientacion/progreso},
\cd{PUT /api/v1/orientacion/progreso}}
\campo{Rationale}{Traduce un proceso administrativo largo en una lista concreta de pasos, que es
lo que el estudiante necesita para saber qué le falta sin interpretar estados internos.}
\campo{Criterio de aceptación}{}
\begin{itemize}
    \item El catálogo de pasos es cerrado y ordenado: tema definido, tutor asignado,
          anteproyecto elaborado, anteproyecto aprobado, solicitud registrada, correcciones
          aplicadas, documento final y pre-sustentación programada.
    \item El porcentaje es el número de pasos completados sobre el total, redondeado al entero
          más próximo.
    \item Un usuario cuyo rol no es \texttt{ESTUDIANTE} responde \texttt{403}.
    \item La consulta y la actualización operan exclusivamente sobre el progreso del estudiante
          en sesión.
\end{itemize}
\campo{Método de verificación}{Prueba}
\campo{Estado}{\estadoV}
\campo{Nota}{Los ocho pasos se marcan manualmente y no se derivan del estado real de la
solicitud (sección~\ref{sec:estados}), de modo que la ruta y el expediente pueden discrepar. La
sincronización automática queda como deuda declarada (\S\ref{sec:pendientes}).}

\subsection{Catálogos académicos y directorio}

\REQ{RF-52}{Gestión de catálogos académicos}
\campo{Enunciado}{El sistema debe permitir a quien posea el permiso
\cd{CARRERAS\_GESTIONAR} crear, editar y eliminar facultades, carreras, modalidades de
titulación y períodos académicos.}
\campo{Prioridad}{Must}
\campo{Fuente}{Auditoría de cobertura del 2026-09-09. Historia de usuario pendiente.}
\campo{Endpoints}{\cd{POST|PUT|DELETE /api/catalogos/facultades},
\cd{POST|PUT|DELETE /api/catalogos/carreras},
\cd{POST|PUT|DELETE /api/catalogos/modalidades},
\cd{POST|PUT|DELETE /api/catalogos/periodos-academicos}}
\campo{Rationale}{Son los datos maestros de los que cuelgan la solicitud, el estudiante y la
convocatoria. Sin ellos el sistema no puede recibir una sola solicitud.}
\campo{Criterio de aceptación}{}
\begin{itemize}
    \item Sin el permiso \cd{CARRERAS\_GESTIONAR} las doce operaciones responden
          \texttt{403}. La consulta permanece abierta a cualquier usuario autenticado (RF-53).
    \item Un código duplicado dentro del mismo catálogo responde \texttt{400}.
    \item Eliminar un elemento referenciado por solicitudes, estudiantes o convocatorias
          responde \texttt{400} en vez de dejar referencias huérfanas.
    \item Toda alta, modificación y baja queda registrada en la bitácora de auditoría con su
          actor (RNF-18).
\end{itemize}
\campo{Método de verificación}{Prueba}
\campo{Estado}{\estadoV}
\campo{Nota}{El Coordinador \textbf{no} posee este permiso, pese a que la tabla de actores de la
especificación anterior le atribuía la gestión de catálogos. La tabla de la sección~\ref{sec:actores}
queda corregida en consecuencia: la gestión de datos maestros es exclusiva del Administrador
(sección~\ref{sec:asimetrias}, punto 1).}

\REQ{RF-53}{Consulta de catálogos de apoyo}
\campo{Enunciado}{El sistema debe ofrecer a cualquier usuario autenticado la consulta de
modalidades, líneas de investigación, áreas temáticas, convocatorias, la convocatoria activa,
carreras, períodos académicos y facultades.}
\campo{Prioridad}{Should}
\campo{Fuente}{HU-02 (los catálogos alimentan el formulario de solicitud)}
\campo{Endpoints}{\cd{GET /api/catalogos/modalidades},
\cd{GET /api/catalogos/lineas-investigacion}, \cd{GET /api/catalogos/areas-tematicas},
\cd{GET /api/catalogos/convocatorias}, \cd{GET /api/catalogos/convocatoria-activa},
\cd{GET /api/catalogos/carreras}, \cd{GET /api/catalogos/periodos-academicos},
\cd{GET /api/catalogos/facultades}}
\campo{Rationale}{Son los valores admisibles que RF-13 valida. Exponerlos como consulta abierta
al usuario autenticado evita que el cliente los duplique en código y se desincronicen.}
\campo{Criterio de aceptación}{}
\begin{itemize}
    \item Las ocho consultas responden \texttt{200} a cualquier usuario autenticado.
    \item La consulta de convocatoria activa devuelve \texttt{404} si no hay ninguna activa, en
          vez de una lista vacía que pueda confundirse con «no hay convocatorias».
    \item Las áreas temáticas admiten filtrado por línea de investigación.
\end{itemize}
\campo{Método de verificación}{Prueba}
\campo{Estado}{\estadoV}

\REQ{RF-54}{Gestión de estudiantes}
\campo{Enunciado}{El sistema debe permitir a quien posea el permiso
\cd{ESTUDIANTES\_GESTIONAR} registrar y editar el perfil académico de un estudiante
---carrera, semestre, período de ingreso y estado académico--- y consultarlo de forma paginada.}
\campo{Prioridad}{Must}
\campo{Fuente}{Auditoría de cobertura del 2026-09-09. Historia de usuario pendiente.}
\campo{Endpoints}{\cd{GET /api/estudiantes/paginado}, \cd{GET /api/estudiantes/\{id\}},
\cd{POST /api/estudiantes}, \cd{PUT /api/estudiantes/\{id\}},
\cd{GET /api/estudiantes/estados-academicos}}
\campo{Rationale}{El perfil de estudiante es distinto de la cuenta de usuario: una cuenta
identifica a una persona, el perfil la sitúa en una carrera y un período, que es lo que la
solicitud necesita.}
\campo{Criterio de aceptación}{}
\begin{itemize}
    \item Sin el permiso \cd{ESTUDIANTES\_GESTIONAR} las cinco operaciones responden
          \texttt{403}.
    \item El estado académico pertenece al conjunto cerrado \texttt{ACTIVO},
          \texttt{EGRESADO}, \texttt{GRADUADO}, \texttt{RETIRADO}, \texttt{SUSPENDIDO};
          cualquier otro responde \texttt{400}.
    \item Un usuario solo puede tener un perfil de estudiante; un segundo registro responde
          \texttt{400}.
    \item El alta genera el código de expediente del estudiante, único en el sistema.
    \item Toda alta y modificación queda registrada en la bitácora de auditoría (RNF-18).
\end{itemize}
\campo{Método de verificación}{Prueba}
\campo{Estado}{\estadoV}

\REQ{RF-55}{Consulta del directorio de docentes}
\campo{Enunciado}{El sistema debe permitir a cualquier usuario autenticado consultar los
docentes, de forma completa, paginada, por identificador, por usuario asociado y filtrando los
disponibles para asignación.}
\campo{Prioridad}{Should}
\campo{Fuente}{HU-03 (la asignación de tutor y tribunal necesita el directorio)}
\campo{Endpoints}{\cd{GET /api/docentes}, \cd{GET /api/docentes/paginado},
\cd{GET /api/docentes/disponibles}, \cd{GET /api/docentes/\{id\}},
\cd{GET /api/docentes/usuario/\{usuarioId\}}}
\campo{Rationale}{El filtro de disponibles es el que hace utilizable la asignación de RF-22 y
RF-27: sin él, la coordinación descubre la indisponibilidad al recibir el rechazo.}
\campo{Criterio de aceptación}{}
\begin{itemize}
    \item El listado de disponibles excluye a los docentes marcados como no disponibles, que son
          exactamente los que RF-27 rechaza.
    \item Ninguna respuesta incluye la contraseña almacenada del usuario asociado.
    \item Sin autenticación responde \texttt{401}.
\end{itemize}
\campo{Método de verificación}{Prueba}
\campo{Estado}{\estadoV}

\subsection{Reportes, observaciones y auditoría}

\REQ{RF-56}{Reportes de gestión}
\campo{Enunciado}{El sistema debe ofrecer a quien posea el permiso \cd{REPORTES\_VER} los
indicadores de resumen, solicitudes por estado, sustentaciones por período, actas, actividad
docente, resultados por carrera, defensas programadas y estadísticas generales.}
\campo{Prioridad}{Should}
\campo{Fuente}{HU-09, CU-09}
\campo{Endpoints}{\cd{GET /api/reportes/resumen},
\cd{GET /api/reportes/solicitudes-por-estado},
\cd{GET /api/reportes/sustentaciones-por-periodo}, \cd{GET /api/reportes/actas},
\cd{GET /api/reportes/actividad-docente}, \cd{GET /api/reportes/por-carrera},
\cd{GET /api/reportes/defensas}, \cd{GET /api/reportes/estadisticas/json}}
\campo{Rationale}{Es la vista agregada con la que la coordinación decide. Los reportes cruzan
calificaciones y actividad docente nominal, así que su control de acceso importa tanto como su
contenido.}
\campo{Criterio de aceptación}{}
\begin{itemize}
    \item Las ocho consultas responden \texttt{200} para \texttt{ADMIN} y \texttt{COORDINADOR},
          y \texttt{403} para \texttt{DOCENTE} y \texttt{ESTUDIANTE}.
    \item Las cifras agregadas coinciden con el conteo por estado de RF-18 sobre los mismos
          datos.
    \item Los estados usados en los agregados son exclusivamente los de la
          sección~\ref{sec:estados}.
\end{itemize}
\campo{Método de verificación}{Prueba + Demostración}
\campo{Estado}{\estadoV}

\REQ{RF-57}{Emisión de reportes en PDF}
\campo{Enunciado}{El sistema debe emitir en formato PDF el reporte de cronograma y el reporte de
estadísticas, a quien posea el permiso \cd{REPORTES\_VER}.}
\campo{Prioridad}{Should}
\campo{Fuente}{HU-09, CU-09}
\campo{Endpoints}{\cd{GET /api/reportes/cronograma/pdf},
\cd{GET /api/reportes/estadisticas/pdf}}
\campo{Rationale}{Los reportes se archivan y se firman fuera del sistema; un PDF es el formato
que la coordinación puede adjuntar a un expediente físico.}
\campo{Criterio de aceptación}{}
\begin{itemize}
    \item Responde \texttt{200} con tipo de contenido \texttt{application/pdf} y el nombre de
          archivo en la cabecera de disposición de contenido.
    \item Sin el permiso \cd{REPORTES\_VER} responde \texttt{403}.
    \item El PDF refleja los mismos datos que el reporte equivalente de RF-56 en el mismo
          instante.
\end{itemize}
\campo{Método de verificación}{Prueba + Demostración}
\campo{Estado}{\estadoV}

\REQ{RF-58}{Estadísticas de acompañamiento por tutor}
\campo{Enunciado}{El sistema debe ofrecer las estadísticas de acompañamiento de los tutores
---estudiantes asignados y estado de sus tutorías--- a quien esté habilitado para ello.}
\campo{Prioridad}{Could}
\campo{Fuente}{Auditoría de cobertura del 2026-09-09; ADR-006. Historia de usuario pendiente.}
\campo{Endpoint}{\cd{GET /api/tutores/estadisticas}}
\campo{Rationale}{Permite a la coordinación detectar tutorías detenidas antes de que bloqueen la
conformación del tribunal, que es donde el retraso se vuelve visible y ya es tarde.}
\campo{Criterio de aceptación}{}
\begin{itemize}
    \item Responde \texttt{200} con el número de estudiantes asignados y el estado de las
          tutorías por tutor.
    \item Sin el permiso exigido responde \texttt{403}.
\end{itemize}
\campo{Método de verificación}{Prueba}
\campo{Estado}{\estadoI}
\campo{Nota}{Este endpoint se protege hoy con \cd{EVALUACION\_RUBRICA\_REGISTRAR}, un
permiso de la categoría Evaluación que no describe esta capacidad y que además poseen los
docentes. La asimetría es real y queda como deuda declarada: el permiso adecuado sería
\cd{REPORTES\_VER} o uno propio de tutoría (\S\ref{sec:pendientes}).}

\REQ{RF-59}{Consulta de observaciones de la solicitud}
\campo{Enunciado}{El sistema debe permitir consultar las observaciones registradas sobre una
solicitud a los actores con derecho sobre ese expediente.}
\campo{Prioridad}{Could}
\campo{Fuente}{HU-02. Historia de usuario pendiente para el endpoint dedicado.}
\campo{Endpoint}{\cd{GET /api/observaciones/solicitud/\{solicitudId\}}}
\campo{Rationale}{Reúne en un solo lugar las observaciones que RF-15, RF-21 y RF-25 producen en
momentos distintos del proceso.}
\campo{Criterio de aceptación}{}
\begin{itemize}
    \item Consultar las observaciones de una solicitud ajena responde \texttt{403}.
    \item Una solicitud sin observaciones responde \texttt{200} con una colección vacía.
\end{itemize}
\campo{Método de verificación}{Prueba}
\campo{Estado}{\estadoI}

\REQ{RF-60}{Consulta de la bitácora de auditoría}
\campo{Enunciado}{El sistema debe permitir a quien posea el permiso \cd{AUDITORIA\_VER}
consultar de forma paginada la bitácora de acciones registradas y la lista de tablas auditadas.}
\campo{Prioridad}{Must}
\campo{Fuente}{Criterio de auditoría de la asignatura; OWASP A09}
\campo{Endpoints}{\cd{GET /api/auditoria/paginado}, \cd{GET /api/auditoria/tablas}}
\campo{Rationale}{Registrar acciones y poder consultarlas son capacidades distintas con controles
de acceso distintos. RNF-18 obliga a registrar; este requisito obliga a que exista una vía de
consulta y a que esté restringida, porque la bitácora contiene datos personales de todos los
usuarios.}
\campo{Criterio de aceptación}{}
\begin{itemize}
    \item Sin el permiso \cd{AUDITORIA\_VER} responde \texttt{403}; el permiso es exclusivo
          del Administrador (sección~\ref{sec:asimetrias}, punto 2).
    \item La consulta paginada acepta filtro por tabla y por rango de fechas.
    \item Ninguna entrada devuelta expone la contraseña almacenada de un usuario, ni siquiera en
          el registro de cambios de la tabla de usuarios.
    \item La bitácora no ofrece ninguna operación de modificación ni de borrado por API.
\end{itemize}
\campo{Método de verificación}{Inspección}
\campo{Estado}{\estadoI}
\campo{Depende de}{RNF-18, RNF-19}
\campo{Nota}{No existe clase de prueba dedicada al módulo de auditoría. Los cuatro criterios se
comprobaron por inspección de la migración que crea la bitácora y su disparador; se declara
\estadoI{} y no \estadoV{} porque la matriz no respalda una prueba automatizada.}

\subsection{Continuidad: respaldo y recuperación}

\REQ{RF-61}{Respaldo de la base de datos bajo demanda y programado}
\campo{Enunciado}{El sistema debe permitir a quien posea el permiso
\cd{BACKUPS\_GESTIONAR} generar respaldos completos bajo demanda, generar respaldos
diferenciales, ejecutarlos automáticamente según un cronograma configurable, y listarlos y
descargarlos.}
\campo{Prioridad}{Must}
\campo{Fuente}{\cd{docs/despliegue/BACKUP.md}; criterio de continuidad de la asignatura.
Historia de usuario pendiente.}
\campo{Endpoints}{\cd{GET /api/backups}, \cd{POST /api/backups},
\cd{POST /api/backups/diferencial}, \cd{GET /api/backups/\{nombre\}/descargar},
\cd{DELETE /api/backups/\{nombre\}}, \cd{GET /api/backups/estado},
\cd{POST /api/backups/bases}, \cd{DELETE /api/backups/bases/\{nombre\}}}
\campo{Rationale}{Un respaldo que nadie puede generar ni descargar no protege nada. El respaldo
es además el único mecanismo del sistema capaz de recuperar un expediente académico perdido, de
modo que su prioridad no es negociable.}
\campo{Criterio de aceptación}{}
\begin{itemize}
    \item Sin el permiso \cd{BACKUPS\_GESTIONAR} las ocho operaciones responden
          \texttt{403}; el permiso es exclusivo del Administrador.
    \item Si el respaldo termina sin error pero produce un archivo vacío, la operación responde
          \texttt{500} declarándolo; \textbf{nunca se registra como respaldo válido}.
    \item Un nombre de respaldo que no corresponda al formato esperado responde \texttt{400}, y
          en ningún caso permite acceder a una ruta fuera del directorio de respaldos.
    \item El listado informa por cada respaldo su tipo, su origen, su fecha y su tamaño.
\end{itemize}
\campo{Método de verificación}{Demostración + Inspección}
\campo{Estado}{\estadoI}
\campo{Depende de}{RNF-24}

\REQ{RF-62}{Configuración del cronograma y de la retención de respaldos}
\campo{Enunciado}{El sistema debe permitir consultar y modificar el cronograma de respaldos, su
activación y su política de retención, y aplicar la retención bajo demanda.}
\campo{Prioridad}{Should}
\campo{Fuente}{\cd{docs/despliegue/BACKUP.md}. Historia de usuario pendiente.}
\campo{Endpoints}{\cd{GET /api/backups/config}, \cd{PUT /api/backups/config},
\cd{POST /api/backups/retencion}}
\campo{Rationale}{La retención es lo que separa un directorio que crece sin límite de una
política de continuidad. Fijarla como configuración, y no en código, permite ajustarla al
crecimiento real sin desplegar.}
\campo{Criterio de aceptación}{}
\begin{itemize}
    \item La configuración por omisión es: respaldo completo los domingos a las 23:00; retención
          de \textbf{7} diarios, \textbf{5} semanales y \textbf{12} mensuales; \textbf{14} días
          de archivado WAL; diferencial desactivado.
    \item Una expresión de cronograma inválida responde \texttt{400} nombrando el valor recibido
          y el formato esperado, y \textbf{no} se guarda ninguna parte del cambio.
    \item La consulta de configuración informa el próximo disparo calculado y el objetivo de
          punto de recuperación que ese cronograma implica.
    \item Todo cambio de configuración registra qué usuario lo realizó y cuándo.
\end{itemize}
\campo{Método de verificación}{Demostración + Inspección}
\campo{Estado}{\estadoI}
\campo{Depende de}{RNF-24}

\REQ{RF-63}{Restauración y prueba de restauración}
\campo{Enunciado}{El sistema debe permitir restaurar un respaldo y ejecutar una prueba de
restauración que verifique que el respaldo es recuperable, registrando su resultado.}
\campo{Prioridad}{Must}
\campo{Fuente}{\cd{docs/despliegue/BACKUP.md}. Historia de usuario pendiente.}
\campo{Endpoints}{\cd{POST /api/backups/\{nombre\}/restaurar},
\cd{GET /api/backups/pruebas}, \cd{POST /api/backups/pruebas}}
\campo{Rationale}{Un respaldo no probado es una suposición. La prueba periódica es la única
forma de saber que el respaldo servirá el día que haga falta, y su registro es la evidencia de
que se hizo.}
\campo{Criterio de aceptación}{}
\begin{itemize}
    \item La prueba de restauración registra su resultado como \texttt{EXITOSA} o
          \texttt{FALLIDA}, con fecha, y ese registro es consultable.
    \item La prueba de restauración \textbf{no} altera la base de datos en servicio.
    \item Sin el permiso \cd{BACKUPS\_GESTIONAR} las tres operaciones responden
          \texttt{403}.
    \item Una restauración fallida deja el sistema en un estado consistente y declara el motivo,
          en vez de dejar la base parcialmente sobrescrita.
\end{itemize}
\campo{Método de verificación}{Demostración}
\campo{Estado}{\estadoI}
\campo{Depende de}{RNF-24}

\REQ{RF-64}{Archivado WAL y recuperación a un punto en el tiempo}
\campo{Enunciado}{El sistema debe permitir consultar el estado del archivado del registro de
escritura anticipada, forzar su rotación y depurar los segmentos que superen el período de
retención configurado.}
\campo{Prioridad}{Should}
\campo{Fuente}{\cd{docs/despliegue/BACKUP.md}. Historia de usuario pendiente.}
\campo{Endpoints}{\cd{GET /api/backups/wal}, \cd{POST /api/backups/wal/switch},
\cd{POST /api/backups/wal/limpiar}}
\campo{Rationale}{Es lo que hace posible recuperar a un instante distinto del último respaldo
completo, y por tanto lo que determina el objetivo de punto de recuperación real de RNF-24. Sin
depuración, el archivado crece hasta agotar el disco.}
\campo{Criterio de aceptación}{}
\begin{itemize}
    \item La consulta informa el número de segmentos archivados, su tamaño total y el segmento
          más antiguo disponible.
    \item La depuración elimina exclusivamente los segmentos anteriores al período de retención
          configurado en RF-62, y nunca los necesarios para recuperar el respaldo completo más
          reciente.
    \item Sin el permiso \cd{BACKUPS\_GESTIONAR} las tres operaciones responden
          \texttt{403}.
\end{itemize}
\campo{Método de verificación}{Demostración}
\campo{Estado}{\estadoI}
\campo{Depende de}{RNF-24}

% ================== 8. REQUISITOS NO FUNCIONALES =============================
\section{Requisitos no funcionales}
\label{sec:rnf}

Los requisitos de esta sección usan la misma plantilla que los funcionales. Su fuente es, en la
mayoría de los casos, una decisión de arquitectura o una norma externa y no una historia de
usuario; se declara explícitamente cuál en vez de dejar el campo vacío.

\subsection{Rendimiento y eficiencia}

\REQ{RNF-01}{Tiempo de respuesta de la API}
\campo{Enunciado}{El sistema debe responder las peticiones REST con un percentil 95 de latencia
inferior a 500\,ms bajo una carga de 50 usuarios concurrentes.}
\campo{Prioridad}{Must} \quad \campo{Fuente}{Criterio de rendimiento de la asignatura}
\campo{Rationale}{Fija un umbral exigible y no una aspiración: sin umbral, ningún requisito se
incumpliría aunque la latencia subiera un orden de magnitud.}
\campo{Criterio de aceptación}{El percentil 95 de \cd{http\_req\_duration} debe ser inferior
a 500\,ms en cada una de al menos 5 corridas independientes de carga con 50 usuarios virtuales
concurrentes, con los resúmenes crudos archivados en el repositorio.}
\campo{Método de verificación}{Análisis}
\campo{Estado}{\estadoV}
\campo{Evidencia}{5 corridas válidas con p95 de 7,70 / 7,39 / 7,47 / 8,53 / 6,17\,ms; media
\textbf{7,45\,ms} (\cd{k6/README.md}, corridas 3 a 7; las dos últimas del 2026-08-29). Dos
corridas anteriores se conservan versionadas y \textbf{marcadas como inválidas}, y no entran en
la estadística.}

\REQ{RNF-02}{Rendimiento percibido del cliente web}
\campo{Enunciado}{El sistema debe alcanzar una puntuación de rendimiento igual o superior a 80
sobre 100 en la auditoría automatizada del cliente web, en perfil de escritorio y de móvil.}
\campo{Prioridad}{Should} \quad \campo{Fuente}{Criterio de calidad de frontend de la asignatura}
\campo{Rationale}{RNF-01 mide el servicio, no lo que el usuario percibe. Un servicio de 7\,ms
puede seguir sintiéndose lento si el cliente tarda seis segundos en pintar.}
\campo{Criterio de aceptación}{Puntuación de rendimiento $\geq 80$ en las corridas de escritorio
y de móvil sobre el build de producción, con los informes archivados y fechados.}
\campo{Método de verificación}{Análisis}
\campo{Estado}{\estadoN}
\campo{Evidencia}{Medición del 2026-09-06, 6 corridas: escritorio \textbf{68--69}, móvil
\textbf{61}. Causa raíz identificada y verificada: tiempo de ejecución de JavaScript bajo la
limitación de CPU $4\times$ del perfil móvil, no peso de recursos.}
\campo{Nota}{Se declara \estadoN{} y no «pendiente de evidencia»: la medición existe y el
resultado no alcanza el umbral. Bajar el umbral para que cuadre sería falsear el requisito. El
plan de cierre está en \S\ref{sec:pendientes}.}

\REQ{RNF-03}{Caché de lectura con vigencia acotada}
\campo{Enunciado}{El sistema debe cachear en memoria distribuida las lecturas de alto tráfico y
la respuesta del proveedor externo, con un tiempo de vigencia declarado por cada conjunto de
datos.}
\campo{Prioridad}{Should} \quad \campo{Fuente}{ADR-001}
\campo{Rationale}{La diferencia medida entre caché fría y caliente en el endpoint externo es de
dos órdenes de magnitud; la vigencia acotada es lo que impide servir datos obsoletos a cambio de
esa ganancia.}
\campo{Criterio de aceptación}{Vigencia por omisión de \textbf{5 minutos}; \textbf{10 minutos}
para el directorio externo de RF-45 y \textbf{5 minutos} para el listado de solicitudes.
Vencida la vigencia, la siguiente lectura debe consultar el origen.}
\campo{Método de verificación}{Análisis + Inspección}
\campo{Estado}{\estadoV}
\campo{Evidencia}{Caché fría $n=30$: media 109,62\,ms, IC 95\,\% [102,57; 116,67]. Caché
caliente $n=30$: IC 95\,\% [7,58; 8,14]. Prueba de Wilcoxon con $U=0$ sobre 900 comparaciones
posibles: separación total (\cd{k6/README.md}).}

\REQ{RNF-04}{Degradación explícita ante indisponibilidad de la caché distribuida}
\campo{Enunciado}{Cuando el almacén de sesión distribuido no esté disponible, el sistema debe
degradar de forma explícita y documentada en cada servicio que depende de él, y en particular
debe \textbf{rechazar} toda petición cuya verificación de revocación de token no pueda
realizarse.}
\campo{Prioridad}{Must} \quad \campo{Fuente}{ADR-002; OWASP A07}
\campo{Rationale}{La verificación de revocación (RF-03) y el límite de tasa (RNF-09) se apoyan
en ese almacén. Si su caída se resuelve dejando pasar la petición, cerrar sesión deja de
revocar nada y el límite de tasa desaparece exactamente cuando la infraestructura está
degradada, que es el peor momento posible. La política es una decisión de seguridad, no un
detalle de infraestructura, y por eso es un requisito y no una nota de riesgo.}
\campo{Criterio de aceptación}{}
\begin{itemize}
    \item \textbf{Revocación de tokens:} con el almacén caído, una petición autenticada responde
          \texttt{401}, no \texttt{200}.
    \item \textbf{Límite de tasa de autenticación:} con el almacén caído, el intento de
          autenticación responde \texttt{503} declarando la degradación, no \texttt{500} ni
          acceso sin límite.
    \item \textbf{Caché de lectura:} con el almacén caído, la lectura se resuelve contra el
          origen y la respuesta se sirve igualmente; es el único servicio donde se admite seguir
          adelante.
    \item Cada degradación deja una entrada de registro que la identifica.
\end{itemize}
\campo{Método de verificación}{Prueba}
\campo{Estado}{\estadoN}
\campo{Nota}{Verificado por inspección el 2026-09-09: la comprobación de revocación captura
cualquier excepción del almacén y devuelve «no revocado», es decir, hoy la política real es
\textbf{fail-open} y un token revocado vuelve a aceptarse mientras el almacén esté caído. El
límite de tasa, en cambio, propaga la excepción y produce \texttt{500}. Ninguno de los dos
comportamientos estaba especificado. Este requisito fija el comportamiento exigido y declara
que hoy no se cumple, en vez de describir lo construido.}

\subsection{Seguridad}

\REQ{RNF-05}{Almacenamiento de contraseñas}
\campo{Enunciado}{El sistema debe almacenar las contraseñas mediante una función de derivación
de clave con salto aleatorio y factor de coste igual o superior a 10, y no debe devolver nunca
la contraseña almacenada en ninguna respuesta.}
\campo{Prioridad}{Must} \quad \campo{Fuente}{OWASP A02; ADR-004}
\campo{Rationale}{Es la única defensa que queda en pie si la base de datos se filtra.}
\campo{Criterio de aceptación}{Ningún registro de usuario contiene la contraseña en claro; dos
usuarios con la misma contraseña producen representaciones almacenadas distintas; ningún cuerpo
de respuesta de RF-01, RF-04, RF-08 ni RF-55 incluye ese campo.}
\campo{Método de verificación}{Prueba + Inspección}
\campo{Estado}{\estadoV}

\REQ{RNF-06}{Política de contraseñas}
\campo{Enunciado}{El sistema debe exigir a toda contraseña una longitud mínima de 8 caracteres y
rechazar las que figuren en una lista de contraseñas comunes, aplicando la misma exigencia en el
alta, en el cambio y en el restablecimiento.}
\campo{Prioridad}{Must} \quad \campo{Fuente}{OWASP A07; NIST SP 800-63B}
\campo{Rationale}{RNF-05 especifica cómo se guarda la contraseña, no qué se exige de ella.
Almacenar impecablemente una contraseña de seis caracteres no protege la cuenta.}
\campo{Criterio de aceptación}{Una contraseña de menos de 8 caracteres responde \texttt{400} en
el alta (RF-04), en el cambio (RF-06) y en el restablecimiento (RF-05); la respuesta declara la
regla incumplida sin revelar la contraseña.}
\campo{Método de verificación}{Prueba}
\campo{Estado}{\estadoN}
\campo{Nota}{La exigencia real vigente es de \textbf{6} caracteres y se aplica \textbf{solo en
el alta}: la actualización de usuario no toca la contraseña y no existen el cambio propio ni el
restablecimiento. Se declara el umbral exigido y el incumplimiento, no el valor actual como si
fuera el requisito.}

\REQ{RNF-07}{Vigencia de los tokens de sesión}
\campo{Enunciado}{El token de acceso debe tener una vigencia de 1 hora y el token de renovación
una vigencia de 7 días, ambas configurables por variable de entorno sin recompilar.}
\campo{Prioridad}{Must} \quad \campo{Fuente}{ADR-002}
\campo{Rationale}{Estos dos valores son normativos: acotan la ventana de uso de una credencial
robada. Hasta esta versión estaban únicamente en un archivo de configuración, es decir, eran
valores normativos escondidos fuera de la especificación.}
\campo{Criterio de aceptación}{El token de acceso expira a los \textbf{3\,600\,000\,ms} y el de
renovación a los \textbf{604\,800\,000\,ms}; un token de acceso expirado responde \texttt{401}
con un mensaje que el cliente pueda distinguir de «sin permiso».}
\campo{Método de verificación}{Prueba + Inspección}
\campo{Estado}{\estadoI}
\campo{Nota}{La vigencia del token de acceso es configurable por entorno; la del token de
renovación está fijada en el archivo de configuración y no se puede ajustar por entorno. El
requisito exige que ambas lo sean.}

\REQ{RNF-08}{Rotación del token de renovación y detección de reutilización}
\campo{Enunciado}{Cada uso de un token de renovación debe invalidarlo y emitir uno nuevo, y la
presentación de un token ya utilizado debe revocar todas las sesiones activas de ese usuario.}
\campo{Prioridad}{Must} \quad \campo{Fuente}{ADR-002; OWASP A07}
\campo{Rationale}{Es lo que convierte el robo de un token de renovación en un evento detectable
en vez de un acceso silencioso durante siete días.}
\campo{Criterio de aceptación}{El criterio observable es el de RF-02; adicionalmente, la
revocación masiva debe dejar sin efecto \textbf{todos} los tokens de renovación del usuario, no
solo el reutilizado, y registrar una alerta identificando al usuario afectado.}
\campo{Método de verificación}{Prueba}
\campo{Estado}{\estadoI}

\REQ{RNF-09}{Límite de tasa en la autenticación}
\campo{Enunciado}{El sistema debe limitar los intentos de autenticación a un máximo de 6 por
minuto y por dirección de origen, respondiendo \texttt{429} al superarlo.}
\campo{Prioridad}{Must} \quad \campo{Fuente}{OWASP A07; ADR-004}
\campo{Rationale}{Sin límite, la política de contraseñas de RNF-06 se puede vencer por fuerza
bruta en línea con independencia de su robustez.}
\campo{Criterio de aceptación}{El séptimo intento dentro de una ventana de \textbf{60 segundos}
desde la misma dirección responde \texttt{429}; el contador expira solo al cerrarse la ventana;
detrás de un proxy inverso se considera la dirección de origen real declarada por el proxy y no
la del proxy.}
\campo{Método de verificación}{Prueba + Análisis}
\campo{Estado}{\estadoV}
\campo{Evidencia}{Corrida dedicada archivada en \cd{k6/rate-limiter-evidence-run.json}.}
\campo{Depende de}{RNF-04}

\REQ{RNF-10}{Cabeceras de seguridad HTTP}
\campo{Enunciado}{Toda respuesta del sistema debe incluir las cabeceras de seguridad que
restringen el origen de los recursos ejecutables, prohíben el enmarcado de la aplicación,
impiden la interpretación especulativa del tipo de contenido y exigen transporte cifrado.}
\campo{Prioridad}{Must} \quad \campo{Fuente}{OWASP A05; ADR-004}
\campo{Rationale}{Son la defensa en profundidad que sigue actuando cuando falla una validación
de entrada.}
\campo{Criterio de aceptación}{}
\begin{itemize}
    \item \texttt{Strict-Transport-Security} con vigencia de \textbf{31\,536\,000\,s} e
          inclusión de subdominios.
    \item \texttt{Content-Security-Policy} con \texttt{script-src 'self'}, \textbf{sin}
          \texttt{'unsafe-inline'} ni \texttt{'unsafe-eval'}, \texttt{connect-src 'self'} y
          \texttt{frame-ancestors 'none'}.
    \item \texttt{X-Frame-Options: DENY} y \texttt{X-Content-Type-Options: nosniff}.
    \item Las mismas cabeceras se emiten en las \textbf{dos} capas que responden al navegador:
          el proxy inverso y el servicio de aplicación.
\end{itemize}
\campo{Método de verificación}{Análisis + Inspección}
\campo{Estado}{\estadoV}
\campo{Evidencia}{Escaneo dinámico del 2026-08-29: \texttt{FAIL-NEW: 0}, \texttt{WARN-NEW: 3},
\texttt{PASS: 58}, \textbf{0 alertas de riesgo alto}. Verificación con petición de cabeceras
transcrita en \cd{docs/mediciones/sec/owasp/OWASP-AUDIT.md}.}

\REQ{RNF-11}{Transporte cifrado}
\campo{Enunciado}{En despliegue accesible desde fuera de la red local, todo tráfico debe viajar
sobre TLS 1.2 o superior, con certificado válido emitido por una autoridad reconocida, y las
peticiones en claro deben redirigirse a su equivalente cifrada.}
\campo{Prioridad}{Must} \quad \campo{Fuente}{OWASP A02; ADR-007}
\campo{Rationale}{La cabecera de transporte estricto de RNF-10 solo tiene efecto si existe
realmente un extremo cifrado; declararla sin TLS es una promesa vacía.}
\campo{Criterio de aceptación}{Sobre el despliegue público: la conexión negocia TLS 1.2 o
superior, el certificado valida contra la cadena pública, una petición en claro responde con
redirección permanente a la versión cifrada, y la cabecera de transporte estricto se emite en la
respuesta cifrada.}
\campo{Método de verificación}{Demostración sobre el despliegue público}
\campo{Estado}{\estadoP}
\campo{Nota}{El sistema \textbf{no está desplegado públicamente} en el punto de corte, de modo
que este requisito no puede verificarse todavía. El proveedor elegido y el procedimiento están
en \cd{docs/despliegue/DEPLOYMENT.md}. La terminación TLS ocurrirá en el borde del proveedor, no
en el proxy inverso del sistema; ese reparto debe quedar declarado aquí cuando el despliegue
exista. No se declara una URL ni un certificado que no se hayan verificado funcionando.}

\REQ{RNF-12}{Autorización resuelta en el servidor}
\campo{Enunciado}{Toda decisión de autorización debe resolverse en el servidor contra el
catálogo de permisos del rol del usuario en sesión, y \textbf{nunca} a partir de un rol,
permiso o identificador enviado por el cliente.}
\campo{Prioridad}{Must} \quad \campo{Fuente}{OWASP A01; ADR-005}
\campo{Rationale}{Es la propiedad que sostiene toda la matriz de la sección~\ref{sec:permisos}.
Si el rol pudiera llegar en el cuerpo de la petición, esa matriz sería decorativa.}
\campo{Criterio de aceptación}{Una petición que declare en su cuerpo o en sus cabeceras un rol o
un permiso distinto del que el usuario en sesión posee obtiene exactamente la misma respuesta
que si no lo declarara; una petición sin autenticación a un endpoint protegido responde
\texttt{401} y una autenticada sin permiso responde \texttt{403}, y ambos casos son
distinguibles por el cliente.}
\campo{Método de verificación}{Prueba + Inspección}
\campo{Estado}{\estadoV}

\REQ{RNF-13}{Aislamiento de los datos por titular}
\campo{Enunciado}{El sistema debe rechazar todo acceso a un recurso que no pertenezca al usuario
en sesión ni corresponda a un rol funcional que este ejerza sobre ese recurso, aunque se conozca
su identificador.}
\campo{Prioridad}{Must} \quad \campo{Fuente}{OWASP A01; 3 hallazgos reales de la auditoría del
2026-08-11/12}
\campo{Rationale}{Es el control que impide que conocer un número de expediente equivalga a poder
leerlo. Se especifica una vez y se referencia desde cada requisito afectado, en vez de repetir la
regla en veinte criterios de aceptación.}
\campo{Criterio de aceptación}{Para cada uno de los recursos solicitud, anteproyecto, fase de
tutoría, evaluación, acta, historial de acta, notificación y perfil, una petición de un usuario
sin derecho sobre ese recurso responde \texttt{403} conociendo su identificador, y la respuesta
no revela si el recurso existe.}
\campo{Método de verificación}{Prueba + Demostración}
\campo{Estado}{\estadoV}

\REQ{RNF-14}{Ausencia de secretos por omisión}
\campo{Enunciado}{El sistema no debe arrancar si el secreto de firma de tokens no está
provisto por el entorno, y no debe existir ningún valor por omisión para ese secreto.}
\campo{Prioridad}{Must} \quad \campo{Fuente}{OWASP A02; hallazgo real de la auditoría del
2026-08-29}
\campo{Rationale}{Un valor por omisión convierte el secreto en público: la auditoría encontró
que el anterior era un secreto de ejemplo que circula en tutoriales. Fallar al arrancar es
preferible a arrancar con una firma que cualquiera puede reproducir.}
\campo{Criterio de aceptación}{Sin la variable de entorno del secreto, el arranque falla con un
error explícito y el servicio no queda a la escucha; el secreto no aparece en ningún registro ni
en ninguna respuesta.}
\campo{Método de verificación}{Inspección + Demostración}
\campo{Estado}{\estadoV}

\REQ{RNF-15}{Ausencia de cuentas de demostración en despliegue no local}
\campo{Enunciado}{El sistema no debe sembrar cuentas con contraseña predefinida fuera del perfil
de desarrollo local, y ninguna contraseña debe figurar literalmente en el código fuente.}
\campo{Prioridad}{Must} \quad \campo{Fuente}{OWASP A05 y A07}
\campo{Rationale}{Una cuenta de administración con contraseña conocida es equivalente a no tener
autenticación. El repositorio es público, de modo que cualquier contraseña escrita en el código
está publicada.}
\campo{Criterio de aceptación}{}
\begin{itemize}
    \item La siembra de datos de demostración se ejecuta \textbf{únicamente} bajo un perfil de
          desarrollo declarado, y no en el arranque por omisión.
    \item Ninguna contraseña aparece literalmente en el código fuente; las cuentas iniciales de
          un despliegue real toman su contraseña del entorno.
    \item Un despliegue no local sin cuenta de administración provista falla al arrancar en vez
          de crear una con contraseña conocida.
\end{itemize}
\campo{Método de verificación}{Inspección}
\campo{Estado}{\estadoN}
\campo{Nota}{Verificado por inspección el 2026-09-09: la siembra se ejecuta en \textbf{todo}
arranque, sin condición de perfil, y crea cuentas de Administrador, Coordinador, Docente y
Estudiante con contraseñas escritas literalmente en el código fuente de un repositorio público.
El archivo \cd{README.md} declara que las credenciales de demostración no se publican «para
evitar dejarlas expuestas en el historial público»; esa afirmación no se sostiene mientras las
contraseñas estén en el código. Es el hallazgo de seguridad más serio del punto de corte y
encabeza el plan de cierre (\S\ref{sec:pendientes}).}

\REQ{RNF-16}{Ausencia de construcción dinámica de consultas}
\campo{Enunciado}{El sistema no debe construir sentencias SQL por concatenación de valores
provenientes de la entrada del usuario.}
\campo{Prioridad}{Must} \quad \campo{Fuente}{OWASP A03; ADR-006}
\campo{Rationale}{La estrategia híbrida de ADR-006 reparte la lógica entre el mapeador
objeto-relacional y procedimientos almacenados, y ninguno de los dos caminos admite concatenar
entrada del usuario. Verificarlo con una herramienta en cada integración, y no por revisión
manual, es lo que impide que una consulta construida a mano se cuele en una entrega futura.}
\campo{Criterio de aceptación}{El análisis estático del backend, con la regla de detección de
SQL dinámico habilitada, reporta \textbf{0} hallazgos de esa categoría en cada corrida
archivada.}
\campo{Método de verificación}{Análisis}
\campo{Estado}{\estadoV}
\campo{Evidencia}{Análisis estático del 2026-08-29: 233 hallazgos totales, \textbf{0} de SQL
dinámico (\cd{docs/mediciones/sec/static-analysis/STATIC-ANALYSIS.md}).}

\REQ{RNF-17}{Gestión de componentes de terceros vulnerables}
\campo{Enunciado}{El sistema no debe entregar al navegador ningún componente de terceros con
vulnerabilidad conocida de severidad alta o crítica, y debe auditar sus dependencias en cada
integración continua.}
\campo{Prioridad}{Should} \quad \campo{Fuente}{OWASP A06}
\campo{Rationale}{Distingue lo que llega al navegador de lo que solo se usa para construir el
proyecto: son superficies de ataque distintas y confundirlas lleva a inflar o a minimizar el
riesgo.}
\campo{Criterio de aceptación}{Cero vulnerabilidades altas o críticas en dependencias
\textbf{enviadas al navegador}; las que residan exclusivamente en herramientas de construcción
se declaran, con su recuento y su justificación, en el informe de auditoría.}
\campo{Método de verificación}{Análisis}
\campo{Estado}{\estadoV}
\campo{Evidencia}{Auditoría del 2026-08-29: la biblioteca de cliente vulnerable se actualizó y
la alerta de riesgo alto pasó a conforme; quedan \textbf{14} vulnerabilidades (desde 41), todas
en herramientas de construcción de iconos. El backend no está auditado contra base de datos de
vulnerabilidades, y así se declara.}

\subsection{Trazabilidad y protección de datos}

\REQ{RNF-18}{Registro de auditoría de las acciones}
\campo{Enunciado}{El sistema debe registrar, para cada alta, modificación y baja sobre las
entidades auditadas, la tabla, el registro afectado, la acción, el usuario que la realizó y el
instante, y ese registro no debe ser modificable ni eliminable desde la aplicación.}
\campo{Prioridad}{Must} \quad \campo{Fuente}{OWASP A09; criterio de auditoría de la asignatura}
\campo{Rationale}{Una bitácora que la aplicación puede editar no es evidencia. Que la escriba el
motor de base de datos, y no la capa de aplicación, es lo que la hace resistente a un fallo o a
un atajo del código.}
\campo{Criterio de aceptación}{}
\begin{itemize}
    \item Cada operación auditada produce exactamente una entrada con actor identificado.
    \item Cuando la operación no proviene de una sesión de usuario, la entrada se registra
          igualmente con el actor sin identificar, en vez de omitirse.
    \item La entrada correspondiente a la tabla de usuarios \textbf{excluye} el campo de
          contraseña, tanto en los datos anteriores como en los nuevos.
    \item No existe ninguna operación de la API que permita modificar o borrar una entrada.
\end{itemize}
\campo{Método de verificación}{Inspección + Prueba}
\campo{Estado}{\estadoV}

\REQ{RNF-19}{Retención y supresión de datos personales}
\campo{Enunciado}{El sistema debe declarar el período de conservación de cada categoría de dato
personal que trata, depurar automáticamente lo que exceda ese período, y ofrecer un mecanismo de
supresión a solicitud del titular compatible con la conservación del expediente académico.}
\campo{Prioridad}{Must} \quad \campo{Fuente}{Ley Orgánica de Protección de Datos Personales del
Ecuador; \cd{docs/etica/ETHICS.md}}
\campo{Rationale}{El sistema trata nombre, correo institucional, teléfono, expediente académico
nominal, calificaciones y una bitácora de acciones por usuario. Hoy nada de eso tiene plazo de
conservación ni vía de supresión: la bitácora crece indefinidamente y el titular no tiene forma
de ejercer sus derechos. Un hallazgo ético sin requisito no tiene fecha de cierre.}
\campo{Criterio de aceptación}{}
\begin{itemize}
    \item Existe una tabla publicada que declara, por categoría de dato, su finalidad, su base
          legal y su período de conservación.
    \item Las entradas de bitácora anteriores al período declarado se depuran de forma
          automática y verificable.
    \item Existe un procedimiento de supresión a solicitud del titular que distingue el dato
          suprimible del que la normativa académica obliga a conservar, y deja constancia de la
          solicitud y de su resolución.
    \item Los formularios de consentimiento firmados se conservan al menos 2 años, conforme al
          protocolo ético del proyecto.
\end{itemize}
\campo{Método de verificación}{Inspección + Demostración}
\campo{Estado}{\estadoP}
\campo{Nota}{Ninguno de los cuatro criterios está construido. El protocolo ético y la plantilla
de consentimiento existen y son una base sólida; lo que falta es convertirlos en obligación con
criterio verificable y fecha de cierre.}

\REQ{RNF-20}{Minimización de datos hacia terceros}
\campo{Enunciado}{El sistema no debe transmitir ningún dato personal de sus usuarios a servicios
externos, incluidos proveedores de modelos de lenguaje.}
\campo{Prioridad}{Must} \quad \campo{Fuente}{Ley Orgánica de Protección de Datos Personales;
\cd{docs/etica/ETHICS.md}}
\campo{Rationale}{Hoy esta propiedad se cumple porque el asistente virtual (RF-44) y la
generación de ideas (RF-47) se resuelven dentro del sistema, y porque el único proveedor externo
(RF-45) recibe una consulta fija sin parámetros del usuario. Pero eso vive solo en el código:
sin requisito, cualquier refactorización futura podría empezar a enviar mensajes de usuarios a
un tercero sin que nada lo impida.}
\campo{Criterio de aceptación}{}
\begin{itemize}
    \item La única petición saliente a un tercero es la consulta del directorio de universidades
          de RF-45, y su contenido no depende de ningún dato del usuario.
    \item El mensaje que el usuario envía al asistente virtual no sale del sistema.
    \item Incorporar un proveedor externo nuevo exige declararlo previamente en la
          sección~\ref{sec:iface-consumidas} y especificar qué datos concretos viajan.
\end{itemize}
\campo{Método de verificación}{Inspección}
\campo{Estado}{\estadoV}

\subsection{Mantenibilidad, usabilidad y operación}

\REQ{RNF-21}{Cobertura de pruebas automatizadas}
\campo{Enunciado}{El sistema debe mantener una cobertura de pruebas igual o superior al 70\,\%
de líneas y de ramas sobre los paquetes de controladores, servicios y seguridad.}
\campo{Prioridad}{Must} \quad \campo{Fuente}{Criterio de mantenibilidad de la asignatura}
\campo{Rationale}{El umbral se declara sobre los paquetes donde vive la lógica, y se excluyen
configuración, entidades y objetos de transferencia: incluirlos inflaría la cifra sin mejorar la
garantía.}
\campo{Criterio de aceptación}{Cobertura $\geq 70$\,\% de líneas y $\geq 70$\,\% de ramas en los
paquetes medidos, con el informe crudo archivado y fechado, regenerado en cada integración
continua.}
\campo{Método de verificación}{Análisis}
\campo{Estado}{\estadoV}
\campo{Evidencia}{Medición del \textbf{2026-09-06}: \textbf{81,03\,\%} de líneas
(3\,148/3\,885) y \textbf{64,39\,\%} de ramas (1\,016/1\,578) en el conjunto; controladores
72,00\,\%/77,41\,\% y servicios 87,57\,\%/70,04\,\%; 559 pruebas en 48 clases.}
\campo{Nota}{La cobertura de ramas del conjunto (64,39\,\%) queda por debajo del 70\,\% aunque
cada uno de los dos paquetes principales lo supere; la diferencia proviene de los paquetes de
seguridad y enumeraciones. Se declara en vez de reportar solo la cifra de líneas. La versión
1.0.0 de este documento declaraba 38,88\,\% y «no cumplido»: esa cifra era de una medición del
2026-08-30 que ya estaba obsoleta cuando se publicó.}

\REQ{RNF-22}{Usabilidad medida}
\campo{Enunciado}{El sistema debe alcanzar una puntuación media igual o superior a 68 sobre 100
en la escala estándar de usabilidad, aplicada a una muestra declarada de usuarios reales de los
cuatro roles.}
\campo{Prioridad}{Should} \quad \campo{Fuente}{Brooke (1996); \cd{docs/etica/ETHICS.md}}
\campo{Rationale}{68 es la media poblacional de referencia del instrumento, de modo que el
umbral significa «al menos tan usable como un sistema medio» y no un número elegido para que
cuadre.}
\campo{Criterio de aceptación}{Media $\geq 68$ sobre una muestra de al menos 10 participantes
reales, con el tamaño de muestra, la composición por rol y los datos crudos anonimizados
archivados, y consentimiento informado firmado por cada participante.}
\campo{Método de verificación}{Análisis}
\campo{Estado}{\estadoP}
\campo{Nota}{El instrumento está listo y no se ha aplicado a personas reales. Una versión
anterior del proyecto publicó una media de usabilidad con diez participantes que \textbf{nunca
existieron}; esos datos fueron retirados y el hecho queda declarado aquí y en
\cd{docs/mediciones/sus/SUS-RESULTS.md}. Se prefiere un requisito sin verificar a una cifra
fabricada.}

\REQ{RNF-23}{Accesibilidad}
\campo{Enunciado}{El cliente web debe alcanzar una puntuación de accesibilidad igual o superior
a 90 sobre 100 en la auditoría automatizada, en perfil de escritorio y de móvil, tomando como
referencia el nivel AA de las pautas WCAG 2.1.}
\campo{Prioridad}{Must} \quad \campo{Fuente}{WCAG 2.1 AA; criterio de calidad de la asignatura}
\campo{Rationale}{El sistema lo usa toda la comunidad universitaria de la carrera; la
accesibilidad no es opcional. La auditoría automatizada no agota WCAG, y por eso el nivel de
referencia se declara además del umbral numérico.}
\campo{Criterio de aceptación}{Puntuación de accesibilidad $\geq 90$ en las corridas de
escritorio y de móvil sobre el build de producción, con los informes archivados y fechados.}
\campo{Método de verificación}{Análisis}
\campo{Estado}{\estadoV}
\campo{Evidencia}{Medición del 2026-09-06, 6 corridas: \textbf{100/100} en escritorio y en
móvil. La corrida del 2026-08-17 daba 89, por debajo del umbral, y se corrigió.}

\REQ{RNF-24}{Objetivos de respaldo y recuperación}
\campo{Enunciado}{El sistema debe garantizar un objetivo de punto de recuperación no superior a
24 horas y un objetivo de tiempo de recuperación no superior a 4 horas, y debe verificar la
recuperabilidad de sus respaldos con una prueba de restauración al menos mensual.}
\campo{Prioridad}{Must} \quad \campo{Fuente}{\cd{docs/despliegue/BACKUP.md}; criterio de
continuidad de la asignatura}
\campo{Rationale}{Sin objetivos declarados, la funcionalidad de respaldo de RF-61 a RF-64 no
tiene contra qué medirse: no se puede saber si una configuración es suficiente. El archivado del
registro de escritura anticipada es lo que permite bajar el objetivo de punto de recuperación
por debajo del intervalo entre respaldos completos.}
\campo{Criterio de aceptación}{}
\begin{itemize}
    \item El cronograma configurado en RF-62 arroja un objetivo de punto de recuperación
          declarado y no superior a \textbf{24 horas}.
    \item Existe al menos una prueba de restauración con resultado \texttt{EXITOSA} en los
          últimos 30 días, registrada conforme a RF-63.
    \item El tiempo medido de una restauración completa sobre datos de tamaño representativo es
          inferior a \textbf{4 horas}, con la medición archivada y fechada.
    \item El período de retención del archivado de escritura anticipada cubre al menos el
          intervalo entre dos respaldos completos consecutivos.
\end{itemize}
\campo{Método de verificación}{Demostración + Análisis}
\campo{Estado}{\estadoI}
\campo{Nota}{La configuración por omisión (respaldo completo semanal) implica un objetivo de
punto de recuperación de hasta 7 días, incompatible con el criterio de 24 horas salvo que el
archivado de escritura anticipada de RF-64 esté activo y verificado. Ese vínculo debe
comprobarse y medirse antes de declarar el requisito verificado.}

\REQ{RNF-25}{Reproducibilidad del entorno y determinismo de la siembra}
\campo{Enunciado}{Desde una clonación limpia del repositorio, el sistema debe poder construirse,
desplegarse, migrarse, probarse y auditarse con una sola orden que termine con código de salida
0, y el estado inicial de los catálogos del dominio debe ser idéntico en dos despliegues
independientes.}
\campo{Prioridad}{Should} \quad \campo{Fuente}{ADR-007; criterio de reproducibilidad de la
asignatura}
\campo{Rationale}{La primera mitad está resuelta y verificada; la segunda no. Los estados de la
solicitud y del cronograma se crean bajo demanda (nota de la sección~\ref{sec:estados}), de modo
que dos despliegues que ejecuten los casos de uso en orden distinto obtienen catálogos con
identificadores distintos. Un sistema que se construye igual pero arranca distinto no es
reproducible.}
\campo{Criterio de aceptación}{}
\begin{itemize}
    \item La orden única de reproducibilidad termina con código 0 desde una clonación limpia.
    \item Las imágenes de contenedor están ancladas por resumen criptográfico, no por etiqueta
          móvil.
    \item Los nueve estados de la solicitud y el estado del cronograma se \textbf{siembran por
          migración}, con identificador y orden fijos; ninguno se crea bajo demanda.
\end{itemize}
\campo{Método de verificación}{Demostración + Inspección}
\campo{Estado}{\estadoN}
\campo{Evidencia}{La orden única está verificada de punta a punta con código de salida 0
(construcción, contenedores, migraciones, 559 pruebas, carga, auditoría, reportes y compilación
del informe), y las imágenes están ancladas por resumen. El tercer criterio no se cumple.}

\REQ{RNF-26}{Versionado uniforme de la API}
\campo{Enunciado}{Todos los endpoints de la API deben exponerse bajo un prefijo de versión
único y explícito.}
\campo{Prioridad}{Should} \quad \campo{Fuente}{ADR-001}
\campo{Rationale}{Un prefijo aplicado a una parte de la API no versiona nada y además induce
error en la documentación: 10 de los 15 endpoints que citaba la matriz anterior no
resolvían precisamente porque se les había añadido un prefijo de versión que su controlador no
tiene.}
\campo{Criterio de aceptación}{El 100\,\% de las rutas de la API responde bajo el prefijo de
versión declarado; una petición a la ruta sin versionar responde \texttt{404} o redirige de
forma permanente a la versionada.}
\campo{Método de verificación}{Inspección}
\campo{Estado}{\estadoN}
\campo{Nota}{En el punto de corte, \textbf{8 de 31} controladores exponen sus rutas bajo el
prefijo de versión y los \textbf{23} restantes no, pese a que el archivo \cd{README.md} declara
la API como versionada. Es la causa directa del defecto de trazabilidad más extendido de la
especificación anterior (\S\ref{sec:cambios}).}

% ================= 9. TRAZABILIDAD Y VERIFICACIÓN ============================
\section{Trazabilidad y verificación}
\label{sec:trazabilidad}

\subsection{Cláusula de verificación}

Ninguna versión de este documento se publica si el validador de trazabilidad no termina con
código de salida 0. El validador es la garantía automática de que este SRS, la matriz y el árbol
de archivos describen el mismo sistema; la revisión manual no puede sostener esa coherencia
sobre 90 requisitos y 207 rutas.

\begin{table}[H]\centering\small
\begin{tabular}{p{0.8cm} p{13.9cm}}
\toprule
\textbf{\#} & \textbf{Comprobación que el validador debe realizar} \\
\midrule
V1 & Toda fila de la matriz tiene exactamente el número de columnas de la cabecera, con los
campos que contienen comas correctamente entrecomillados. \\
V2 & La columna de estado contiene únicamente los cuatro valores del vocabulario cerrado de la
sección~\ref{sec:estados-req}; cualquier otro texto falla. \\
V3 & Todo identificador de requisito de este SRS existe en la matriz, y todo identificador de la
matriz existe en este SRS. El conteo debe coincidir. \\
V4 & Todo endpoint citado en la matriz resuelve contra un controlador real: se comprueba la
\textbf{ruta completa}, método y camino, no solo su primer segmento. \\
V5 & Toda clase de prueba citada existe como archivo en disco. \\
V6 & Todo archivo de historia de usuario o de caso de uso citado existe como archivo en disco. \\
V7 & Ningún requisito declarado \emph{Verificado} carece de prueba automatizada o de evidencia
archivada. \\
V8 & Ningún requisito declarado \emph{Planificado} cita un endpoint existente: si el endpoint
existe, el requisito está al menos implementado. \\
\bottomrule
\end{tabular}
\caption{Las ocho comprobaciones exigidas. V1 a V6 son nuevas y son
exactamente las que habrían detectado, antes de la entrega, los defectos que la
sección~\ref{sec:cambios} documenta.}
\end{table}

\subsection{Estructura de la matriz de trazabilidad}

La matriz (\cd{docs/trazabilidad/matriz.csv}) conecta cada requisito con su origen y su
evidencia. Sus columnas son:

\begin{table}[H]\centering\small
\begin{tabular}{p{4.0cm} p{10.7cm}}
\toprule
\textbf{Columna} & \textbf{Contenido} \\
\midrule
\cd{ID\_Requisito} & Identificador \texttt{RF-nn} o \texttt{RNF-nn} de este documento. \\
\texttt{Titulo} & Título del requisito, idéntico al de este documento. \\
\texttt{Fuente} & Historia de usuario, caso de uso, decisión de arquitectura o norma. \\
\texttt{Modulo} & Módulo funcional del sistema. \\
\cd{Endpoint\_REST} & Método y ruta completa, tal como responde el servicio. \\
\cd{Prioridad\_MoSCoW} & \emph{Must}, \emph{Should} o \emph{Could}. \\
\cd{Metodo\_Verificacion} & Prueba, Demostración, Análisis o Inspección. \\
\cd{Caso\_Prueba} & Clase o clases de prueba, o \texttt{ninguna}. \\
\texttt{Estado} & Uno de los cuatro valores del vocabulario cerrado, y ningún otro. \\
\texttt{Evidencia} & Ruta del artefacto de evidencia con su fecha. \\
\texttt{Observaciones} & Todo matiz que no cabe en el vocabulario cerrado de estado. \\
\bottomrule
\end{tabular}
\end{table}

\noindent La columna \texttt{Observaciones} es nueva y existe por una razón concreta: es donde
va el matiz que antes se colaba dentro de la columna de estado y rompía el vocabulario.

% ================= 10. ESTADO DE CUMPLIMIENTO ================================
\section{Estado de cumplimiento del corpus}
\label{sec:cumplimiento}

\subsection{Distribución por estado y prioridad}

\begin{table}[H]\centering\small
\begin{tabular}{lrrrrr}
\toprule
& \textbf{Verificado} & \textbf{Implementado} & \textbf{Planificado} & \textbf{No cumplido} &
\textbf{Total} \\
\midrule
Requisitos funcionales     & 46 & 16 & 2 & 0 & \textbf{64} \\
Requisitos no funcionales  & 14 &  3 & 3 & 6 & \textbf{26} \\
\midrule
\textbf{Total}             & \textbf{60} & \textbf{19} & \textbf{5} & \textbf{6} &
\textbf{90} \\
\bottomrule
\end{tabular}
\caption{Distribución por estado. Contado sobre la etiqueta \texttt{v1.0.1}.}
\end{table}

\begin{table}[H]\centering\small
\begin{tabular}{lrrrr}
\toprule
& \textbf{Must} & \textbf{Should} & \textbf{Could} & \textbf{Total} \\
\midrule
Requisitos funcionales    & 30 & 26 & 8 & \textbf{64} \\
Requisitos no funcionales & 20 &  6 & 0 & \textbf{26} \\
\midrule
\textbf{Total}            & \textbf{50} & \textbf{32} & \textbf{8} & \textbf{90} \\
\bottomrule
\end{tabular}
\caption{Distribución por prioridad MoSCoW. No hay requisitos \emph{Won't}: una capacidad
descartada se retira del alcance, no se conserva como requisito con etiqueta de descarte.}
\end{table}

\subsection{Requisitos \emph{Must} que no están verificados}

De los \textbf{50} requisitos de prioridad \emph{Must}, \textbf{36 están verificados
(72,0\,\%)}. Los \textbf{14} restantes se enumeran aquí de forma completa, porque un porcentaje
sin el detalle de lo que falta no es información accionable.

\begin{longtable}{p{1.6cm} p{5.3cm} p{2.9cm} p{4.6cm}}
\toprule
\textbf{ID} & \textbf{Requisito} & \textbf{Estado} & \textbf{Qué falta} \\
\midrule
\endfirsthead
\toprule
\textbf{ID} & \textbf{Requisito} & \textbf{Estado} & \textbf{Qué falta} \\
\midrule
\endhead
RNF-15 & Ausencia de cuentas de demostración & \estadoN & Retirar del código las contraseñas
literales y condicionar la siembra a un perfil de desarrollo. \\
RNF-04 & Degradación ante caída de la caché & \estadoN & Cambiar la verificación de revocación a
\emph{fail-closed} y dar respuesta declarada al límite de tasa. \\
RNF-06 & Política de contraseñas & \estadoN & Subir el mínimo a 8 caracteres y aplicarlo en las
tres vías de alta y cambio. \\
RF-05 & Recuperación de contraseña & \estadoP & Construir la funcionalidad completa. \\
RF-06 & Cambio de contraseña propia & \estadoP & Construir la funcionalidad completa. \\
RNF-19 & Retención y supresión de datos & \estadoP & Tabla de retención, depuración automática y
procedimiento de supresión. \\
RNF-11 & Transporte cifrado & \estadoP & Depende del despliegue público, aún no realizado. \\
RNF-07 & Vigencia de tokens & \estadoI & Hacer configurable por entorno la vigencia del token de
renovación. \\
RNF-08 & Rotación y detección de reutilización & \estadoI & Prueba automatizada del escenario de
reutilización. \\
RF-02 & Renovación con rotación & \estadoI & Prueba automatizada del escenario de reutilización. \\
RNF-24 & Objetivos de respaldo y recuperación & \estadoI & Medir el tiempo de restauración y
verificar el vínculo con el archivado de escritura anticipada. \\
RF-61 & Respaldo bajo demanda y programado & \estadoI & Evidencia archivada de una corrida
completa. \\
RF-63 & Restauración y prueba de restauración & \estadoI & Evidencia archivada de una prueba de
restauración exitosa. \\
RF-60 & Consulta de la bitácora de auditoría & \estadoI & Clase de prueba dedicada al módulo de
auditoría. \\
\bottomrule
\caption{Los 14 requisitos \emph{Must} no verificados, con lo que falta en cada uno.}
\end{longtable}

\noindent Ocho de los catorce se cierran con evidencia o con una prueba; seis exigen cambiar
código. Los seis con implicación de seguridad o legal encabezan el plan de la
sección~\ref{sec:pendientes}.

\subsection{Nota metodológica sobre el crecimiento del corpus}

El corpus pasó de los 16 requisitos de la especificación anterior a los 90 de esta. Conviene ser preciso sobre
qué significa esa cifra, porque se presta a dos lecturas opuestas y solo una es correcta.

\textbf{No es alcance nuevo.} De los 74 requisitos añadidos, \textbf{69 documentan capacidades o
propiedades que ya estaban construidas} en el repositorio y que sencillamente no tenían requisito
que las especificara. Solo \textbf{5} describen algo que aún no existe o no puede verificarse
todavía: RF-05 y RF-06 (credenciales del usuario), RNF-19 (retención y supresión), RNF-22
(usabilidad medida) y RNF-11 (transporte cifrado, a la espera del despliegue público). La tasa de
adición no mide crecimiento del producto: mide cuánto le faltaba a la especificación para
alcanzar al código.

\textbf{Es, sin embargo, una señal sobre el proceso.} Que un SRS describiera 16 requisitos de un
sistema con 207 rutas significa que durante la mayor parte del desarrollo el documento dejó de
funcionar como contrato y como base de verificación. En un proyecto con interesados externos,
una desviación de esta magnitud detectada al cierre de una entrega sería una señal de alarma
sobre la especificación continua, no un logro de la revisión final. Se declara aquí y se remite
al capítulo de amenazas a la validez del informe.

\textbf{Estabilidad del corpus heredado.} De los 16 requisitos de la especificación anterior, 15 se
conservan con el mismo alcance (renumerados y completados con la plantilla) y 1 cambia de
prioridad con justificación (RF-37, de \emph{Should} a \emph{Must}). Ninguno se retira. La
estabilidad del corpus heredado es del \textbf{93,8\,\%}; no debe sumarse ni compararse con la
tasa de adición, porque miden cosas distintas.

% =============== 11. CORRESPONDENCIA CON EL ANEXO C ==========================
\section{Correspondencia con el Anexo C de ISO/IEC/IEEE 29148:2018}
\label{sec:anexoc}

Este documento declara conformidad con la norma, de modo que debe poder contrastarse con el
esquema que la norma propone. La estructura elegida no es idéntica a la del Anexo C; esta tabla
mapea cada elemento del esquema normativo a la sección donde vive, y justifica cada desviación
en vez de dejarla implícita.

\begin{longtable}{p{5.1cm} p{4.1cm} p{5.2cm}}
\toprule
\textbf{Elemento del Anexo C} & \textbf{Sección de este SRS} & \textbf{Desviación y
justificación} \\
\midrule
\endfirsthead
\toprule
\textbf{Elemento del Anexo C} & \textbf{Sección de este SRS} & \textbf{Desviación y
justificación} \\
\midrule
\endhead
Propósito & \S2.1 & Sin desviación. \\
Alcance & \S2.2 & Sin desviación. Se añade el alcance excluido de forma explícita. \\
Definiciones y acrónimos & \S2.3 & Sin desviación. \\
Referencias & \S2.7 & Sin desviación. \\
Perspectiva del producto & \S3.1 & Sin desviación. \\
Funciones del producto & \S7 & Se antepone el catálogo de estados (\S5) y la matriz de permisos
(\S6) porque los criterios de aceptación de \S7 los usan como vocabulario. \\
Características de los usuarios & \S3.2 & Se distingue rol de cuenta de rol funcional, distinción
que la norma no exige pero sin la cual el control de acceso de RF-37 no es especificable. \\
Restricciones & \S3.3 & Separadas de los requisitos, conforme a la norma. \\
Supuestos y dependencias & \S3.4 & Sin desviación. \\
Interfaces externas & \S4 & Se separa interfaz expuesta de interfaz consumida, y se exige a cada
consumida un comportamiento declarado ante indisponibilidad. \\
Requisitos funcionales & \S7 & Sin desviación. \\
Requisitos de rendimiento & \S8.1 & Agrupados como no funcionales, con umbral y condición de
carga. \\
Requisitos lógicos de base de datos & \S5 & \textbf{Desviación:} no hay una sección de modelo
lógico. El ciclo de vida de los datos se especifica como máquinas de estado, que es lo que los
criterios de aceptación necesitan; el modelo físico vive en las migraciones versionadas y su
catálogo de rutinas, referenciados en \S2.7. \\
Restricciones de diseño & \S3.3 & Sin desviación. \\
Atributos del sistema: fiabilidad y disponibilidad & \S8.1 (RNF-04), \S8.4 (RNF-24, RNF-25) &
Sin desviación. \\
Atributos del sistema: seguridad & \S8.2 & Sin desviación. Se mapea cada requisito a su control
de referencia externo. \\
Atributos del sistema: mantenibilidad y portabilidad & \S8.4 & Portabilidad cubierta por RD-04 y
RNF-25 en vez de por requisitos propios. \\
Usabilidad & \S8.4 (RNF-22, RNF-23) & \textbf{Desviación:} usabilidad y accesibilidad se
especifican como requisitos con umbral medible en vez de como sección propia, para que estén
sujetas al mismo vocabulario de estado que el resto. \\
Verificación & \S9.1 & Se declara como cláusula ejecutable (validador con código de salida 0), no
solo como enunciado de intención. \\
Trazabilidad & \S9.2 y \cd{matriz.csv} & Sin desviación. \\
Información de soporte & \S14, \S15 & Anexos de INCOSE y elicitación por referencia, para que las
dos copias no puedan divergir. \\
\bottomrule
\caption{Correspondencia con el esquema del Anexo C. Cuatro desviaciones declaradas, cada una
con su justificación.}
\end{longtable}

% ================ 12. CAMBIOS RESPECTO A LA VERSIÓN 1.0.0 ====================
\section{Correcciones respecto a la especificación anterior}
\label{sec:cambios}

Esta sección documenta lo que se corrigió y por qué, para que la bitácora de cambios no sea la
única constancia. El detalle por requisito está en \cd{docs/requisitos/CHANGELOG-REQ.md}.

\subsection{Defectos de trazabilidad corregidos}

\begin{enumerate}
    \item \textbf{Diez de los quince endpoints citados en la matriz no resolvían.} Todos
          llevaban un prefijo de versión que su controlador no tiene: la matriz citaba, por
          ejemplo, la ruta versionada de la creación de solicitud, del cronograma, de la
          evaluación, del acta, de los reportes, de las salas, de los usuarios y del
          anteproyecto, y ninguna de ellas existe. La causa raíz es RNF-26. Corregido: los 207
          endpoints de este documento se citan con la ruta real.
    \item \textbf{Tres historias de usuario citadas no existían como archivo.} La matriz
          referenciaba HU-13, HU-14 y HU-15; el directorio contiene HU-01 a HU-12. Corregido:
          ningún requisito de esta versión cita una fuente inexistente, y los que aún no tienen
          historia lo declaran explícitamente (\S\ref{sec:pendientes}).
    \item \textbf{El SRS y la matriz declaraban conjuntos distintos.} El documento hablaba de 12
          requisitos y la matriz contenía 15 filas, en la misma versión. Corregido: los dos
          artefactos contienen los mismos 90 identificadores, y la comprobación V3 lo verifica en
          integración continua.
    \item \textbf{El validador no detectaba ninguno de los tres problemas anteriores.}
          Comprobaba el endpoint quedándose con el \textbf{primer segmento} de la ruta tras
          descartar el prefijo, de modo que una ruta versionada inexistente coincidía con el
          controlador sin versionar y pasaba sin protestar. Ese era el único punto ciego de un
          mecanismo que por lo demás funcionaba. Corregido con las comprobaciones V1 a V6 y V8
          de la sección~\ref{sec:trazabilidad}.
\end{enumerate}

\subsection{Estados declarados que ya no correspondían}

\begin{enumerate}
    \item \textbf{La cobertura de pruebas.} La especificación anterior declaraba 38,88\,\% de líneas y
          «no cumplido» contra un umbral de 60\,\%. La medición vigente en ese mismo momento ya
          era otra, y la del punto de corte es de \textbf{81,03\,\%} de líneas con 559 pruebas.
          El requisito pasa a \estadoV{} con umbral elevado al 70\,\% de líneas y de ramas
          (RNF-21). La cifra ahora va fechada dentro del propio requisito.
    \item \textbf{La accesibilidad.} No existía requisito que la exigiera pese a medirse. Se
          especifica con umbral y nivel de referencia (RNF-23), y se declara verificada con
          100/100.
    \item \textbf{El rendimiento del cliente web.} Tampoco existía requisito. Se especifica
          (RNF-02) y se declara \estadoN{} con la cifra real, en vez de omitir la medición que
          no alcanza el umbral.
\end{enumerate}

\subsection{Cobertura: de 16 requisitos a 90}

La especificación anterior describía 16 requisitos de un sistema con 31 controladores y 207 rutas.
Módulos completos no tenían ninguna línea que los describiera: respaldo y recuperación, bitácora
de auditoría, catálogos académicos, gestión de estudiantes, centro de orientación, tutoría por
fases, asistente virtual, estado en tiempo real, roles y permisos, y buena parte del ciclo de
vida de la solicitud y del acta. Esta versión los especifica. Ninguno exigió programar: solo
especificar lo que ya estaba programado.

\subsection{Secciones estructurales incorporadas}

Cinco secciones que la norma espera y el documento no tenía: interfaces externas (\S4), catálogo
de estados del dominio (\S5), matriz de permisos rol\,$\times$\,operación (\S6), correspondencia
con el Anexo C (\S\ref{sec:anexoc}) y cláusula de verificación (\S9.1). Las dos que más falta
hacían son \S5 y \S6: los criterios de aceptación usaban nombres de estado nunca definidos, y
con 29 permisos repartidos entre 207 rutas la autorización ya no se podía auditar leyendo
requisito por requisito.

\subsection{Lenguaje y forma}

Se fija la convención de redacción normativa (\S2.4), se cierra el vocabulario de estado
(\S2.5) y de método de verificación (\S2.6), se expresan los criterios de aceptación como
respuesta observable en vez de como excepción interna, y se retiran del enunciado las rutas de
archivo y los nombres de clase, que pasan a la matriz donde el validador los comprueba contra el
disco.

% ==================== 13. DEUDA DECLARADA Y PLAN =============================
\section{Deuda declarada y plan de cierre}
\label{sec:pendientes}

Nada de lo que sigue está resuelto. Se enumera de forma completa porque un hallazgo sin registro
no tiene fecha de cierre, y porque declararlo es preferible a que aparezca en una revisión
externa.

\subsection{Orden de trabajo propuesto}

\begin{enumerate}
    \item \textbf{Seguridad inmediata.} RNF-15 (contraseñas literales en el código de un
          repositorio público y siembra de cuentas en todo arranque) y RNF-04 (la revocación de
          tokens es hoy \emph{fail-open}). Son los dos únicos hallazgos que degradan la seguridad
          del sistema desplegado, y ambos se cierran con un cambio acotado.
    \item \textbf{Credenciales del usuario.} RF-05, RF-06 y RNF-06: recuperación, cambio y
          política de contraseñas. Hoy no existe ninguna vía, ni propia ni administrativa, para
          cambiar una contraseña, y el asistente virtual anuncia una capacidad que el sistema no
          tiene.
    \item \textbf{Protección de datos personales.} RNF-19: tabla de retención por categoría,
          depuración automática de la bitácora y procedimiento de supresión a solicitud del
          titular. El protocolo ético del proyecto es una base sólida; falta convertirlo en
          requisito con criterio y fecha.
    \item \textbf{Comprobaciones del validador.} Implementar V1 a V6 y V8
          (\S\ref{sec:trazabilidad}). Son pocas líneas y son lo que impide que esta revisión se
          repita: las seis habrían detectado solas los cuatro defectos de \S\ref{sec:cambios}.
    \item \textbf{Continuidad.} RNF-24, RF-61 y RF-63: medir el tiempo de restauración, archivar
          la evidencia de una prueba exitosa y verificar que el archivado de escritura anticipada
          sostiene el objetivo de punto de recuperación declarado.
    \item \textbf{Despliegue y transporte.} RNF-11, que hoy no puede verificarse porque el
          sistema no está desplegado públicamente, y RNF-02, cuyo plan de cierre es la
          fragmentación del paquete de JavaScript del cliente.
    \item \textbf{Coherencia interna.} RNF-25 (sembrar por migración los estados de la solicitud
          y del cronograma), RNF-26 (unificar el prefijo de versión de la API) y la revisión
          permiso a permiso de la asignación del Coordinador
          (\S\ref{sec:asimetrias}, punto 4).
\end{enumerate}

\subsection{Historias de usuario pendientes de redactar}

Veinticuatro requisitos de esta versión declaran su fuente como «auditoría de cobertura» porque
documentan módulos construidos sin historia previa. Sin fuente no se puede sostener que un
requisito sea necesario, así que la historia hace falta aunque el código exista: RF-05, RF-06,
RF-10, RF-16, RF-23, RF-24, RF-25, RF-26, RF-38, RF-39, RF-40, RF-43, RF-44, RF-46, RF-47,
RF-48, RF-49, RF-50, RF-51, RF-52, RF-54, RF-58, RF-61 y RF-62 a RF-64. La comprobación V6 del
validador impedirá, en adelante, citar una historia que no exista.

\subsection{Defectos funcionales menores declarados}

\begin{enumerate}
    \item \textbf{Ruta de la \emph{cookie} de renovación.} La \emph{cookie} del token de
          renovación se emite con un ámbito de ruta que no coincide con el del endpoint que la
          consume, de modo que el navegador no la envía en la petición de renovación. Afecta a
          RF-02.
    \item \textbf{Estado \texttt{SUSPENDIDA} sin retorno.} No existe la reanudación de una
          solicitud suspendida (RF-16), de modo que una suspensión es hoy tan definitiva como un
          rechazo.
    \item \textbf{Ruta de progreso desacoplada del expediente.} Los ocho pasos de RF-51 se marcan
          manualmente y pueden contradecir el estado real de la solicitud.
    \item \textbf{Permiso inadecuado en las estadísticas de tutoría.} RF-58 se protege con un
          permiso de la categoría Evaluación.
    \item \textbf{Colisión de identificadores en el catálogo de permisos.} Dos migraciones
          asignaron el mismo identificador a permisos distintos; la inserción posterior se
          descartó en silencio y hubo que repararla con una migración adicional. No hay efecto
          vigente, pero el catálogo no es reproducible por simple lectura de las migraciones en
          orden.
    \item \textbf{Salvaguarda de administración.} RF-12 permite retirar al rol de Administrador
          el permiso de gestionar permisos, dejando el sistema sin ningún usuario capaz de
          revertirlo.
\end{enumerate}

% ======================== 14. ANEXOS =========================================
\section{Anexo: \emph{checklist} INCOSE}

Las 9 características de calidad individuales de INCOSE (Necesario, Apropiado, No ambiguo,
Completo, Singular, Factible, Verificable, Correcto, Conforme) y las 6 del conjunto se aplican al
corpus completo. El \emph{checklist} vive en \cd{docs/checklists/INCOSE-REQUIREMENTS.md} y se
anexa por referencia, no verbatim, para que las dos copias no puedan divergir.

\textbf{Brechas conocidas del corpus de esta versión:} la característica \emph{Correcto} no puede
afirmarse para los 5 requisitos en estado \estadoP{}, que describen comportamiento aún no
construido; y \emph{Verificable como conjunto} está limitada por los 19 requisitos en estado
\estadoI{}, cuyo criterio no tiene todavía prueba ni evidencia archivada. La característica
\emph{Conforme}, que era la brecha más consistente de la especificación anterior por usar formato de
historia de usuario en vez de lenguaje normativo, queda cerrada con la convención de \S2.4.

\section{Anexo: técnicas de elicitación aplicadas}

Detalle completo en \cd{docs/requisitos/elicitacion/README.md}. De las cinco técnicas esperadas,
\textbf{dos tienen evidencia real y verificable} en el repositorio (análisis de documentos, con
observaciones docentes trazadas a \emph{commits} reales, y evaluación cruzada de pares),
\textbf{una es parcial} (prototipado iterativo evidenciado por el historial de control de
versiones, sin sesiones formales documentadas) y \textbf{dos no tienen evidencia documentada}
(entrevistas y observación directa). Se declara explícitamente en vez de fabricar transcripciones
o actas que no existieron.

A estas se suma, en esta versión, una técnica que sí produjo la mayor parte del corpus nuevo: la
\textbf{ingeniería inversa de requisitos} sobre el código, las migraciones y la configuración del
sistema construido. Es una técnica legítima de elicitación cuando la implementación precede a la
especificación, y declararla es más honesto que atribuir a entrevistas inexistentes unos
requisitos que se obtuvieron leyendo el repositorio.

% ==================== 15. PÁGINA DE APROBACIÓN ===============================
\section{Página de aprobación}

\begin{center}
\begin{tabular}{p{7cm} p{7cm}}
\toprule
\textbf{Elaborado por (Equipo de desarrollo)} & \textbf{Revisado y aprobado por
(Docente-Director)} \\
\midrule
& \\[1.4cm]
\rule{6cm}{0.4pt} & \rule{6cm}{0.4pt} \\
Alava Alvarado, Jean Pierre & Ing. Guerrero Ulloa Gleiston Cíceron, Mg. \\
\vspace{0.25cm} & \\
\rule{6cm}{0.4pt} & Fecha de aprobación: \rule{4cm}{0.4pt} \\
Moncayo Loor, Xavier Alejandro & \\
\vspace{0.25cm} & \\
\rule{6cm}{0.4pt} & \\
Zamora Arias, Carla Esthefania & \\
\vspace{0.25cm} & \\
\rule{6cm}{0.4pt} & \\
Barreto Rosado, Heider Dominick & \\
\addlinespace
\rule{6cm}{0.4pt} & \\
Representante del equipo & \\
\bottomrule
\end{tabular}
\end{center}

\vspace{0.4cm}
\noindent\fbox{\parbox{0.98\textwidth}{\textbf{Estado real de la firma (2026-09-09):} la firma
del docente-director \textbf{no se ha obtenido todavía}. Requiere su revisión y aprobación real,
fuera del alcance de lo que el equipo puede completar por sí mismo. No se fabrica una firma ni se
marca esta sección como completa hasta que ocurra realmente, siguiendo la misma política de
integridad aplicada al resto de este documento: donde no hay evidencia, se declara que no la
hay.}}

\end{document}
